\documentclass[si.tex]{subfiles}
\usepackage{array}

\begin{document}


\newcommand{\safegraphic}[4]{%
  \IfFileExists{#2}{%
    \adjustbox{valign=t}{\includegraphics[width=#1]{#2}}%
  }{%
    \adjustbox{valign=t}{%
      \fcolorbox{gray}{white}{%
        \begin{minipage}[c][#3][c]{#1}%
        \centering
        \scriptsize Missing PDF\\
        \scriptsize #4
        \end{minipage}%
      }%
    }%
  }%
}

\clearpage

\section{Robustness Studies}\label{sec:alt-noise-loss-experiments}

\subsection{Efficient-Coding Family}\label{sec:simulations-efficient-coding}
\providecommand{\safegraphic}[4]{%
  \IfFileExists{#2}{%
    \adjustbox{valign=t}{\includegraphics[width=#1]{#2}}%
  }{%
    \adjustbox{valign=t}{%
      \fcolorbox{gray}{white}{%
        \begin{minipage}[c][#3][c]{#1}%
        \centering
        \scriptsize Missing PDF\\
        \scriptsize #4
        \end{minipage}%
      }%
    }%
  }%
}


\newcommand{\efficientcodingrowTwoEight}[4]{%
#2 &
\safegraphic{0.14\textwidth}{../code/Synthetic/figures/CounterfactualModel_VIZ_ByNoiseMagnGauge_Components.py_2345_#4_#1_2_0_10.0_180.pdf}{2.2cm}{gt #2 N#3} &
\safegraphic{0.18\textwidth}{../code/Synthetic/figures/RunSynthetic_FreePrior_CosineLoss_OnSim_VIZ_OnlyModel_OtherNoiseLevels.py_SimulateSynthetic_Parameterized_OtherNoiseLevels_Grid_VarySize.py_180_2_2345_N#3_#4_#1.txt_2_0_10.0_180.pdf}{3.0cm}{fit p=2 #2 N#3} &
\safegraphic{0.10\textwidth}{../code/Synthetic/figures/evaluateCrossValidationResults_Synthetic_Gardelle_NonF.py_SimulateSynthetic_Parameterized_OtherNoiseLevels_Grid_VarySize.py_180_2_2345_N#3_#4_#1.txt_RelativeLF.pdf}{3.0cm}{NLL p=2 #2 N#3} &
\safegraphic{0.18\textwidth}{../code/Synthetic/figures/RunSynthetic_FreePrior_CosineLoss_OnSim_VIZ_OnlyModel_OtherNoiseLevels.py_SimulateSynthetic_Parameterized_OtherNoiseLevels_Grid_VarySize.py_180_8_2345_N#3_#4_#1.txt_8_0_10.0_180.pdf}{3.0cm}{fit p=8 #2 N#3} &
\safegraphic{0.10\textwidth}{../code/Synthetic/figures/evaluateCrossValidationResults_Synthetic_Gardelle_NonF.py_SimulateSynthetic_Parameterized_OtherNoiseLevels_Grid_VarySize.py_180_8_2345_N#3_#4_#1.txt_RelativeLF.pdf}{3.0cm}{NLL p=8 #2 N#3} \\
}

\newcommand{\efficientcodingrowTwo}[4]{%
#2 &
\safegraphic{0.14\textwidth}{../code/Synthetic/figures/CounterfactualModel_VIZ_ByNoiseMagnGauge_Components.py_2345_#4_#1_2_0_10.0_180.pdf}{2.2cm}{gt #2 N#3} &
\safegraphic{0.18\textwidth}{../code/Synthetic/figures/RunSynthetic_FreePrior_CosineLoss_OnSim_VIZ_OnlyModel_OtherNoiseLevels.py_SimulateSynthetic_Parameterized_OtherNoiseLevels_Grid_VarySize.py_180_2_2345_N#3_#4_#1.txt_2_0_10.0_180.pdf}{3.0cm}{fit p=2 #2 N#3} &
\safegraphic{0.10\textwidth}{../code/Synthetic/figures/evaluateCrossValidationResults_Synthetic_Gardelle_NonF.py_SimulateSynthetic_Parameterized_OtherNoiseLevels_Grid_VarySize.py_180_2_2345_N#3_#4_#1.txt_RelativeLF.pdf}{3.0cm}{NLL p=2 #2 N#3} \\
}

\begin{figure}[p]
\centering
\begin{tabular}{>{\centering\arraybackslash}m{0.13\textwidth}c@{\hspace{0.020\textwidth}}cc@{\hspace{0.020\textwidth}}cc}
& \textbf{Ground truth} & \multicolumn{2}{c}{\textbf{Ground truth: \(p=2\)}} & \multicolumn{2}{c}{\textbf{Ground truth: \(p=8\)}} \\
\textbf{Exponent} $\alpha$ & \textbf{Model} & \textbf{Fit at \(p=2\)} & \textbf{\(\Delta\)NLL} & \textbf{Fit at \(p=8\)} & \textbf{\(\Delta\)NLL} \\
\efficientcodingrowTwoEight{STEEPPERIODICa0b5}{$1/2$}{10000}{STEEPPERIODIC}
\efficientcodingrowTwoEight{STEEPPERIODIC}{$1$}{10000}{STEEPPERIODIC}
\efficientcodingrowTwoEight{STEEPPERIODICa2b0}{$2$}{10000}{STEEPPERIODIC}
\end{tabular}
\caption{
Identifiability of models where prior and encoding are linked directly via efficient coding. We simulate data (10K trials, four noise levels) where $\Prior$ is periodic, and the encoding resources are directly proportional to either the prior itself, its square root, or its square.
}
\label{fig:si-alt-efficient-coding-periodic}
\end{figure}

\begin{figure}[p]
\centering
\begin{tabular}{>{\centering\arraybackslash}m{0.13\textwidth}ccccc}
& \textbf{Ground truth} & \multicolumn{2}{c}{\textbf{Ground truth: \(p=2\)}}  \\
\textbf{Exponent} $\alpha$ & \textbf{Model} & \textbf{Fit at \(p=2\)} & \textbf{\(\Delta\)NLL}  \\
\efficientcodingrowTwo{FOURIERa0b5_121}{$1/2$}{10000}{FOURIER_121}
\efficientcodingrowTwo{FOURIERa1b0_121}{$1$}{10000}{FOURIER_121}
\efficientcodingrowTwo{FOURIERa2b0_121}{$2$}{10000}{FOURIER_121}
\end{tabular}

\ 

\ 

\begin{tabular}{>{\centering\arraybackslash}m{0.13\textwidth}ccccc}
& \textbf{Ground truth} & \multicolumn{2}{c}{\textbf{Ground truth: \(p=2\)}}  \\
\textbf{Exponent} $\alpha$ & \textbf{Model} & \textbf{Fit at \(p=2\)} & \textbf{\(\Delta\)NLL}  \\
\efficientcodingrowTwo{FOURIERa0b5_122}{$1/2$}{10000}{FOURIER_122}
\efficientcodingrowTwo{FOURIERa1b0_122}{$1$}{10000}{FOURIER_122}
\efficientcodingrowTwo{FOURIERa2b0_122}{$2$}{10000}{FOURIER_122}
\end{tabular}
\caption{
Identifiability of models where prior and encoding are linked directly via efficient coding. We show two randomly generated models (10K trials, four noise levels, $p=2$).
}
\label{fig:si-alt-efficient-coding-fourier}
\end{figure}





\clearpage

\subsection{Robustness to Large Motor Noise}

In our main simulations, we set the motor noise magnitude based on fits to the behavioral datasets collected by \citet{deGardelle2010AnOI} and \citet{Remington2018LateBI}.
Here, we study the effect of larger motor noise on identifiability, by simulating two additional groups of datasets with larger motor noise levels than those used in the main simulations, with all other parameters matching the simulations from Section~\ref{sec:further-circular}. The noise magnitudes are shown in Figure~\ref{fig:motor-noise-magnitudes}, and example fits are shown in Figures~\ref{fig:si-alt-medium-large-motor-noise} and~\ref{fig:si-alt-large-motor-noise}.
Results show that higher motor noise may lead to more jagged fits to prior or encoding and flatter NLL, but models remain identifiable even at the largest motor noise level (Figure \ref{fig:si-alt-large-motor-noise}).

\begin{figure}
\centering
\includegraphics[width=0.75\textwidth]{motor_noise_magnitudes.pdf}
\caption{Motor noise distributions, as von Mises densities on the orientation space $[0^\circ, 180^\circ]$ used in our circular simulations.
The "Orientation (de Gardelle et al)" version is the same as in the main circular simulations, and taken from our fit to the behavioral data collected by \citet{deGardelle2010AnOI}. 
The Medium Large and Large settings are used for the simulations in Figures~\ref{fig:si-alt-medium-large-motor-noise} and~\ref{fig:si-alt-large-motor-noise}.
}\label{fig:motor-noise-magnitudes}
\end{figure}



\begin{figure}[ht]
\centering
\large\textbf{10K Trials}\\[0.5em]
\begin{tabular}{>{\centering\arraybackslash}m{0.17\textwidth}cc@{\hspace{1.2em}}cc}
& \multicolumn{2}{c}{\textbf{Ground truth: \(p=2\)}} & \multicolumn{2}{c}{\textbf{Ground truth: \(p=8\)}} \\
& Fit at \(p=2\) & \(\Delta\)NLL & Fit at \(p=8\) & \(\Delta\)NLL \\
\safegraphic{0.14\textwidth}{figures/CounterfactualModel_VIZ_Components.py_5_STEEPPERIODIC_STEEPPERIODIC_2_0_10.0_180.pdf}{2.2cm}{gt periodic periodic} &
\safegraphic{0.18\textwidth}{../code/Synthetic/figures/RunSynthetic_FreePrior_CosineLoss_OnSim_VIZ_OnlyModel_OtherNoiseLevels.py_SimulateSynthetic_Parameterized_OtherNoiseLevels_Grid_VarySize_LargeMotorNoise.py_180_2_2345_N10000_STEEPPERIODIC_STEEPPERIODIC.txt_2_0_10.0_180.pdf}{3.0cm}{fit p=2} &
\safegraphic{0.10\textwidth}{../code/Synthetic/figures/evaluateCrossValidationResults_Synthetic_LargeMotorNoise.py_SimulateSynthetic_Parameterized_OtherNoiseLevels_Grid_VarySize_LargeMotorNoise.py_180_2_2345_N10000_STEEPPERIODIC_STEEPPERIODIC.txt_RelativeLF.pdf}{3.0cm}{NLL p=2} &
\safegraphic{0.18\textwidth}{../code/Synthetic/figures/RunSynthetic_FreePrior_CosineLoss_OnSim_VIZ_OnlyModel_OtherNoiseLevels.py_SimulateSynthetic_Parameterized_OtherNoiseLevels_Grid_VarySize_LargeMotorNoise.py_180_8_2345_N10000_STEEPPERIODIC_STEEPPERIODIC.txt_8_0_10.0_180.pdf}{3.0cm}{fit p=8} &
\safegraphic{0.10\textwidth}{../code/Synthetic/figures/evaluateCrossValidationResults_Synthetic_LargeMotorNoise.py_SimulateSynthetic_Parameterized_OtherNoiseLevels_Grid_VarySize_LargeMotorNoise.py_180_8_2345_N10000_STEEPPERIODIC_STEEPPERIODIC.txt_RelativeLF.pdf}{3.0cm}{NLL p=8} \\
\safegraphic{0.14\textwidth}{figures/CounterfactualModel_VIZ_Components.py_5_STEEPPERIODIC_UNIFORM_2_0_10.0_180.pdf}{2.2cm}{gt periodic uniform} &
\safegraphic{0.18\textwidth}{../code/Synthetic/figures/RunSynthetic_FreePrior_CosineLoss_OnSim_VIZ_OnlyModel_OtherNoiseLevels.py_SimulateSynthetic_Parameterized_OtherNoiseLevels_Grid_VarySize_LargeMotorNoise.py_180_2_2345_N10000_STEEPPERIODIC_UNIFORM.txt_2_0_10.0_180.pdf}{3.0cm}{fit p=2} &
\safegraphic{0.10\textwidth}{../code/Synthetic/figures/evaluateCrossValidationResults_Synthetic_LargeMotorNoise.py_SimulateSynthetic_Parameterized_OtherNoiseLevels_Grid_VarySize_LargeMotorNoise.py_180_2_2345_N10000_STEEPPERIODIC_UNIFORM.txt_RelativeLF.pdf}{3.0cm}{NLL p=2} &
\safegraphic{0.18\textwidth}{../code/Synthetic/figures/RunSynthetic_FreePrior_CosineLoss_OnSim_VIZ_OnlyModel_OtherNoiseLevels.py_SimulateSynthetic_Parameterized_OtherNoiseLevels_Grid_VarySize_LargeMotorNoise.py_180_8_2345_N10000_STEEPPERIODIC_UNIFORM.txt_8_0_10.0_180.pdf}{3.0cm}{fit p=8} &
\safegraphic{0.10\textwidth}{../code/Synthetic/figures/evaluateCrossValidationResults_Synthetic_LargeMotorNoise.py_SimulateSynthetic_Parameterized_OtherNoiseLevels_Grid_VarySize_LargeMotorNoise.py_180_8_2345_N10000_STEEPPERIODIC_UNIFORM.txt_RelativeLF.pdf}{3.0cm}{NLL p=8} \\
\safegraphic{0.14\textwidth}{figures/CounterfactualModel_VIZ_Components.py_5_UNIFORM_STEEPPERIODIC_2_0_10.0_180.pdf}{2.2cm}{gt uniform periodic} &
\safegraphic{0.18\textwidth}{../code/Synthetic/figures/RunSynthetic_FreePrior_CosineLoss_OnSim_VIZ_OnlyModel_OtherNoiseLevels.py_SimulateSynthetic_Parameterized_OtherNoiseLevels_Grid_VarySize_LargeMotorNoise.py_180_2_2345_N10000_UNIFORM_STEEPPERIODIC.txt_2_0_10.0_180.pdf}{3.0cm}{fit p=2} &
\safegraphic{0.10\textwidth}{../code/Synthetic/figures/evaluateCrossValidationResults_Synthetic_LargeMotorNoise.py_SimulateSynthetic_Parameterized_OtherNoiseLevels_Grid_VarySize_LargeMotorNoise.py_180_2_2345_N10000_UNIFORM_STEEPPERIODIC.txt_RelativeLF.pdf}{3.0cm}{NLL p=2} &
\safegraphic{0.18\textwidth}{../code/Synthetic/figures/RunSynthetic_FreePrior_CosineLoss_OnSim_VIZ_OnlyModel_OtherNoiseLevels.py_SimulateSynthetic_Parameterized_OtherNoiseLevels_Grid_VarySize_LargeMotorNoise.py_180_8_2345_N10000_UNIFORM_STEEPPERIODIC.txt_8_0_10.0_180.pdf}{3.0cm}{fit p=8} &
\safegraphic{0.10\textwidth}{../code/Synthetic/figures/evaluateCrossValidationResults_Synthetic_LargeMotorNoise.py_SimulateSynthetic_Parameterized_OtherNoiseLevels_Grid_VarySize_LargeMotorNoise.py_180_8_2345_N10000_UNIFORM_STEEPPERIODIC.txt_RelativeLF.pdf}{3.0cm}{NLL p=8} \\
\safegraphic{0.14\textwidth}{figures/CounterfactualModel_VIZ_Components.py_5_STEEPSHIFTED_STEEPPERIODIC_2_0_10.0_180.pdf}{2.2cm}{gt shifted periodic} &
\safegraphic{0.18\textwidth}{../code/Synthetic/figures/RunSynthetic_FreePrior_CosineLoss_OnSim_VIZ_OnlyModel_OtherNoiseLevels.py_SimulateSynthetic_Parameterized_OtherNoiseLevels_Grid_VarySize_LargeMotorNoise.py_180_2_2345_N10000_STEEPSHIFTED_STEEPPERIODIC.txt_2_0_10.0_180.pdf}{3.0cm}{fit p=2} &
\safegraphic{0.10\textwidth}{../code/Synthetic/figures/evaluateCrossValidationResults_Synthetic_LargeMotorNoise.py_SimulateSynthetic_Parameterized_OtherNoiseLevels_Grid_VarySize_LargeMotorNoise.py_180_2_2345_N10000_STEEPSHIFTED_STEEPPERIODIC.txt_RelativeLF.pdf}{3.0cm}{NLL p=2} &
\safegraphic{0.18\textwidth}{../code/Synthetic/figures/RunSynthetic_FreePrior_CosineLoss_OnSim_VIZ_OnlyModel_OtherNoiseLevels.py_SimulateSynthetic_Parameterized_OtherNoiseLevels_Grid_VarySize_LargeMotorNoise.py_180_8_2345_N10000_STEEPSHIFTED_STEEPPERIODIC.txt_8_0_10.0_180.pdf}{3.0cm}{fit p=8} &
\safegraphic{0.10\textwidth}{../code/Synthetic/figures/evaluateCrossValidationResults_Synthetic_LargeMotorNoise.py_SimulateSynthetic_Parameterized_OtherNoiseLevels_Grid_VarySize_LargeMotorNoise.py_180_8_2345_N10000_STEEPSHIFTED_STEEPPERIODIC.txt_RelativeLF.pdf}{3.0cm}{NLL p=8}
\end{tabular}
\\[1.2em]
\large\textbf{40K Trials}\\[0.5em]
\begin{tabular}{>{\centering\arraybackslash}m{0.17\textwidth}cc@{\hspace{1.2em}}cc}
& \multicolumn{2}{c}{\textbf{Ground truth: \(p=2\)}} & \multicolumn{2}{c}{\textbf{Ground truth: \(p=8\)}} \\
& Fit at \(p=2\) & \(\Delta\)NLL & Fit at \(p=8\) & \(\Delta\)NLL \\
\safegraphic{0.14\textwidth}{figures/CounterfactualModel_VIZ_Components.py_5_STEEPPERIODIC_STEEPPERIODIC_2_0_10.0_180.pdf}{2.2cm}{gt periodic periodic} &
\safegraphic{0.18\textwidth}{../code/Synthetic/figures/RunSynthetic_FreePrior_CosineLoss_OnSim_VIZ_OnlyModel_OtherNoiseLevels.py_SimulateSynthetic_Parameterized_OtherNoiseLevels_Grid_VarySize_LargeMotorNoise.py_180_2_2345_N40000_STEEPPERIODIC_STEEPPERIODIC.txt_2_0_10.0_180.pdf}{3.0cm}{fit p=2} &
\safegraphic{0.10\textwidth}{../code/Synthetic/figures/evaluateCrossValidationResults_Synthetic_LargeMotorNoise.py_SimulateSynthetic_Parameterized_OtherNoiseLevels_Grid_VarySize_LargeMotorNoise.py_180_2_2345_N40000_STEEPPERIODIC_STEEPPERIODIC.txt_RelativeLF.pdf}{3.0cm}{NLL p=2} &
\safegraphic{0.18\textwidth}{../code/Synthetic/figures/RunSynthetic_FreePrior_CosineLoss_OnSim_VIZ_OnlyModel_OtherNoiseLevels.py_SimulateSynthetic_Parameterized_OtherNoiseLevels_Grid_VarySize_LargeMotorNoise.py_180_8_2345_N40000_STEEPPERIODIC_STEEPPERIODIC.txt_8_0_10.0_180.pdf}{3.0cm}{fit p=8} &
\safegraphic{0.10\textwidth}{../code/Synthetic/figures/evaluateCrossValidationResults_Synthetic_LargeMotorNoise.py_SimulateSynthetic_Parameterized_OtherNoiseLevels_Grid_VarySize_LargeMotorNoise.py_180_8_2345_N40000_STEEPPERIODIC_STEEPPERIODIC.txt_RelativeLF.pdf}{3.0cm}{NLL p=8} \\
\safegraphic{0.14\textwidth}{figures/CounterfactualModel_VIZ_Components.py_5_STEEPPERIODIC_UNIFORM_2_0_10.0_180.pdf}{2.2cm}{gt periodic uniform} &
\safegraphic{0.18\textwidth}{../code/Synthetic/figures/RunSynthetic_FreePrior_CosineLoss_OnSim_VIZ_OnlyModel_OtherNoiseLevels.py_SimulateSynthetic_Parameterized_OtherNoiseLevels_Grid_VarySize_LargeMotorNoise.py_180_2_2345_N40000_STEEPPERIODIC_UNIFORM.txt_2_0_10.0_180.pdf}{3.0cm}{fit p=2} &
\safegraphic{0.10\textwidth}{../code/Synthetic/figures/evaluateCrossValidationResults_Synthetic_LargeMotorNoise.py_SimulateSynthetic_Parameterized_OtherNoiseLevels_Grid_VarySize_LargeMotorNoise.py_180_2_2345_N40000_STEEPPERIODIC_UNIFORM.txt_RelativeLF.pdf}{3.0cm}{NLL p=2} &
\safegraphic{0.18\textwidth}{../code/Synthetic/figures/RunSynthetic_FreePrior_CosineLoss_OnSim_VIZ_OnlyModel_OtherNoiseLevels.py_SimulateSynthetic_Parameterized_OtherNoiseLevels_Grid_VarySize_LargeMotorNoise.py_180_8_2345_N40000_STEEPPERIODIC_UNIFORM.txt_8_0_10.0_180.pdf}{3.0cm}{fit p=8} &
\safegraphic{0.10\textwidth}{../code/Synthetic/figures/evaluateCrossValidationResults_Synthetic_LargeMotorNoise.py_SimulateSynthetic_Parameterized_OtherNoiseLevels_Grid_VarySize_LargeMotorNoise.py_180_8_2345_N40000_STEEPPERIODIC_UNIFORM.txt_RelativeLF.pdf}{3.0cm}{NLL p=8} \\
\safegraphic{0.14\textwidth}{figures/CounterfactualModel_VIZ_Components.py_5_UNIFORM_STEEPPERIODIC_2_0_10.0_180.pdf}{2.2cm}{gt uniform periodic} &
\safegraphic{0.18\textwidth}{../code/Synthetic/figures/RunSynthetic_FreePrior_CosineLoss_OnSim_VIZ_OnlyModel_OtherNoiseLevels.py_SimulateSynthetic_Parameterized_OtherNoiseLevels_Grid_VarySize_LargeMotorNoise.py_180_2_2345_N40000_UNIFORM_STEEPPERIODIC.txt_2_0_10.0_180.pdf}{3.0cm}{fit p=2} &
\safegraphic{0.10\textwidth}{../code/Synthetic/figures/evaluateCrossValidationResults_Synthetic_LargeMotorNoise.py_SimulateSynthetic_Parameterized_OtherNoiseLevels_Grid_VarySize_LargeMotorNoise.py_180_2_2345_N40000_UNIFORM_STEEPPERIODIC.txt_RelativeLF.pdf}{3.0cm}{NLL p=2} &
\safegraphic{0.18\textwidth}{../code/Synthetic/figures/RunSynthetic_FreePrior_CosineLoss_OnSim_VIZ_OnlyModel_OtherNoiseLevels.py_SimulateSynthetic_Parameterized_OtherNoiseLevels_Grid_VarySize_LargeMotorNoise.py_180_8_2345_N40000_UNIFORM_STEEPPERIODIC.txt_8_0_10.0_180.pdf}{3.0cm}{fit p=8} &
\safegraphic{0.10\textwidth}{../code/Synthetic/figures/evaluateCrossValidationResults_Synthetic_LargeMotorNoise.py_SimulateSynthetic_Parameterized_OtherNoiseLevels_Grid_VarySize_LargeMotorNoise.py_180_8_2345_N40000_UNIFORM_STEEPPERIODIC.txt_RelativeLF.pdf}{3.0cm}{NLL p=8} \\
\safegraphic{0.14\textwidth}{figures/CounterfactualModel_VIZ_Components.py_5_STEEPSHIFTED_STEEPPERIODIC_2_0_10.0_180.pdf}{2.2cm}{gt shifted periodic} &
\safegraphic{0.18\textwidth}{../code/Synthetic/figures/RunSynthetic_FreePrior_CosineLoss_OnSim_VIZ_OnlyModel_OtherNoiseLevels.py_SimulateSynthetic_Parameterized_OtherNoiseLevels_Grid_VarySize_LargeMotorNoise.py_180_2_2345_N40000_STEEPSHIFTED_STEEPPERIODIC.txt_2_0_10.0_180.pdf}{3.0cm}{fit p=2} &
\safegraphic{0.10\textwidth}{../code/Synthetic/figures/evaluateCrossValidationResults_Synthetic_LargeMotorNoise.py_SimulateSynthetic_Parameterized_OtherNoiseLevels_Grid_VarySize_LargeMotorNoise.py_180_2_2345_N40000_STEEPSHIFTED_STEEPPERIODIC.txt_RelativeLF.pdf}{3.0cm}{NLL p=2} &
\safegraphic{0.18\textwidth}{../code/Synthetic/figures/RunSynthetic_FreePrior_CosineLoss_OnSim_VIZ_OnlyModel_OtherNoiseLevels.py_SimulateSynthetic_Parameterized_OtherNoiseLevels_Grid_VarySize_LargeMotorNoise.py_180_8_2345_N40000_STEEPSHIFTED_STEEPPERIODIC.txt_8_0_10.0_180.pdf}{3.0cm}{fit p=8} &
\safegraphic{0.10\textwidth}{../code/Synthetic/figures/evaluateCrossValidationResults_Synthetic_LargeMotorNoise.py_SimulateSynthetic_Parameterized_OtherNoiseLevels_Grid_VarySize_LargeMotorNoise.py_180_8_2345_N40000_STEEPSHIFTED_STEEPPERIODIC.txt_RelativeLF.pdf}{3.0cm}{NLL p=8}
\end{tabular}
\caption{Identifiability on datasets sampled at highly increased motor noise magnitude (``large'' motor noise in Figure \ref{fig:si-alt-large-motor-noise}), on data with 10K or 40K trials, generated at $p=2$ (left column) or $p=8$ (right column).
Except for motor noise, all other parameters are chosen as in the main simulations from Section~\ref{sec:further-circular}.
Prior and encoding are recovered approximately, with an improved fit at 40K trials.
The ground-truth loss function or exponents close to it are always singled out, despite the highly increased motor noise magnitude.}
\label{fig:si-alt-large-motor-noise}
\end{figure}

\begin{comment}
# LargeMotorNoise: fitted-model visualizations at the ground-truth exponent (N10000)
python3 RunSynthetic_FreePrior_CosineLoss_OnSim_VIZ_OnlyModel_OtherNoiseLevels.py 2 0 10.0 180 "SimulateSynthetic_Parameterized_OtherNoiseLevels_Grid_VarySize_LargeMotorNoise.py_180_2_2345_N10000_STEEPPERIODIC_STEEPPERIODIC.txt"
python3 RunSynthetic_FreePrior_CosineLoss_OnSim_VIZ_OnlyModel_OtherNoiseLevels.py 2 0 10.0 180 "SimulateSynthetic_Parameterized_OtherNoiseLevels_Grid_VarySize_LargeMotorNoise.py_180_2_2345_N10000_STEEPPERIODIC_UNIFORM.txt"
python3 RunSynthetic_FreePrior_CosineLoss_OnSim_VIZ_OnlyModel_OtherNoiseLevels.py 2 0 10.0 180 "SimulateSynthetic_Parameterized_OtherNoiseLevels_Grid_VarySize_LargeMotorNoise.py_180_2_2345_N10000_UNIFORM_STEEPPERIODIC.txt"
python3 RunSynthetic_FreePrior_CosineLoss_OnSim_VIZ_OnlyModel_OtherNoiseLevels.py 2 0 10.0 180 "SimulateSynthetic_Parameterized_OtherNoiseLevels_Grid_VarySize_LargeMotorNoise.py_180_2_2345_N10000_STEEPSHIFTED_STEEPPERIODIC.txt"
python3 RunSynthetic_FreePrior_CosineLoss_OnSim_VIZ_OnlyModel_OtherNoiseLevels.py 8 0 10.0 180 "SimulateSynthetic_Parameterized_OtherNoiseLevels_Grid_VarySize_LargeMotorNoise.py_180_8_2345_N10000_STEEPPERIODIC_STEEPPERIODIC.txt"
python3 RunSynthetic_FreePrior_CosineLoss_OnSim_VIZ_OnlyModel_OtherNoiseLevels.py 8 0 10.0 180 "SimulateSynthetic_Parameterized_OtherNoiseLevels_Grid_VarySize_LargeMotorNoise.py_180_8_2345_N10000_STEEPPERIODIC_UNIFORM.txt"
python3 RunSynthetic_FreePrior_CosineLoss_OnSim_VIZ_OnlyModel_OtherNoiseLevels.py 8 0 10.0 180 "SimulateSynthetic_Parameterized_OtherNoiseLevels_Grid_VarySize_LargeMotorNoise.py_180_8_2345_N10000_UNIFORM_STEEPPERIODIC.txt"
python3 RunSynthetic_FreePrior_CosineLoss_OnSim_VIZ_OnlyModel_OtherNoiseLevels.py 8 0 10.0 180 "SimulateSynthetic_Parameterized_OtherNoiseLevels_Grid_VarySize_LargeMotorNoise.py_180_8_2345_N10000_STEEPSHIFTED_STEEPPERIODIC.txt"

# LargeMotorNoise: NLL-by-loss-function plots (N10000)
python3 evaluateCrossValidationResults_Synthetic_LargeMotorNoise.py "SimulateSynthetic_Parameterized_OtherNoiseLevels_Grid_VarySize_LargeMotorNoise.py_180_2_2345_N10000_STEEPPERIODIC_STEEPPERIODIC.txt"
python3 evaluateCrossValidationResults_Synthetic_LargeMotorNoise.py "SimulateSynthetic_Parameterized_OtherNoiseLevels_Grid_VarySize_LargeMotorNoise.py_180_2_2345_N10000_STEEPPERIODIC_UNIFORM.txt"
python3 evaluateCrossValidationResults_Synthetic_LargeMotorNoise.py "SimulateSynthetic_Parameterized_OtherNoiseLevels_Grid_VarySize_LargeMotorNoise.py_180_2_2345_N10000_UNIFORM_STEEPPERIODIC.txt"
python3 evaluateCrossValidationResults_Synthetic_LargeMotorNoise.py "SimulateSynthetic_Parameterized_OtherNoiseLevels_Grid_VarySize_LargeMotorNoise.py_180_2_2345_N10000_STEEPSHIFTED_STEEPPERIODIC.txt"
python3 evaluateCrossValidationResults_Synthetic_LargeMotorNoise.py "SimulateSynthetic_Parameterized_OtherNoiseLevels_Grid_VarySize_LargeMotorNoise.py_180_8_2345_N10000_STEEPPERIODIC_STEEPPERIODIC.txt"
python3 evaluateCrossValidationResults_Synthetic_LargeMotorNoise.py "SimulateSynthetic_Parameterized_OtherNoiseLevels_Grid_VarySize_LargeMotorNoise.py_180_8_2345_N10000_STEEPPERIODIC_UNIFORM.txt"
python3 evaluateCrossValidationResults_Synthetic_LargeMotorNoise.py "SimulateSynthetic_Parameterized_OtherNoiseLevels_Grid_VarySize_LargeMotorNoise.py_180_8_2345_N10000_UNIFORM_STEEPPERIODIC.txt"
python3 evaluateCrossValidationResults_Synthetic_LargeMotorNoise.py "SimulateSynthetic_Parameterized_OtherNoiseLevels_Grid_VarySize_LargeMotorNoise.py_180_8_2345_N10000_STEEPSHIFTED_STEEPPERIODIC.txt"

# LargeMotorNoise: fitted-model visualizations at the ground-truth exponent (N40000)
python3 RunSynthetic_FreePrior_CosineLoss_OnSim_VIZ_OnlyModel_OtherNoiseLevels.py 2 0 10.0 180 "SimulateSynthetic_Parameterized_OtherNoiseLevels_Grid_VarySize_LargeMotorNoise.py_180_2_2345_N40000_STEEPPERIODIC_STEEPPERIODIC.txt"
python3 RunSynthetic_FreePrior_CosineLoss_OnSim_VIZ_OnlyModel_OtherNoiseLevels.py 2 0 10.0 180 "SimulateSynthetic_Parameterized_OtherNoiseLevels_Grid_VarySize_LargeMotorNoise.py_180_2_2345_N40000_STEEPPERIODIC_UNIFORM.txt"
python3 RunSynthetic_FreePrior_CosineLoss_OnSim_VIZ_OnlyModel_OtherNoiseLevels.py 2 0 10.0 180 "SimulateSynthetic_Parameterized_OtherNoiseLevels_Grid_VarySize_LargeMotorNoise.py_180_2_2345_N40000_UNIFORM_STEEPPERIODIC.txt"
python3 RunSynthetic_FreePrior_CosineLoss_OnSim_VIZ_OnlyModel_OtherNoiseLevels.py 2 0 10.0 180 "SimulateSynthetic_Parameterized_OtherNoiseLevels_Grid_VarySize_LargeMotorNoise.py_180_2_2345_N40000_STEEPSHIFTED_STEEPPERIODIC.txt"
python3 RunSynthetic_FreePrior_CosineLoss_OnSim_VIZ_OnlyModel_OtherNoiseLevels.py 8 0 10.0 180 "SimulateSynthetic_Parameterized_OtherNoiseLevels_Grid_VarySize_LargeMotorNoise.py_180_8_2345_N40000_STEEPPERIODIC_STEEPPERIODIC.txt"
python3 RunSynthetic_FreePrior_CosineLoss_OnSim_VIZ_OnlyModel_OtherNoiseLevels.py 8 0 10.0 180 "SimulateSynthetic_Parameterized_OtherNoiseLevels_Grid_VarySize_LargeMotorNoise.py_180_8_2345_N40000_STEEPPERIODIC_UNIFORM.txt"
python3 RunSynthetic_FreePrior_CosineLoss_OnSim_VIZ_OnlyModel_OtherNoiseLevels.py 8 0 10.0 180 "SimulateSynthetic_Parameterized_OtherNoiseLevels_Grid_VarySize_LargeMotorNoise.py_180_8_2345_N40000_UNIFORM_STEEPPERIODIC.txt"
python3 RunSynthetic_FreePrior_CosineLoss_OnSim_VIZ_OnlyModel_OtherNoiseLevels.py 8 0 10.0 180 "SimulateSynthetic_Parameterized_OtherNoiseLevels_Grid_VarySize_LargeMotorNoise.py_180_8_2345_N40000_STEEPSHIFTED_STEEPPERIODIC.txt"

# LargeMotorNoise: NLL-by-loss-function plots (N40000)
python3 evaluateCrossValidationResults_Synthetic_LargeMotorNoise.py "SimulateSynthetic_Parameterized_OtherNoiseLevels_Grid_VarySize_LargeMotorNoise.py_180_2_2345_N40000_STEEPPERIODIC_STEEPPERIODIC.txt"
python3 evaluateCrossValidationResults_Synthetic_LargeMotorNoise.py "SimulateSynthetic_Parameterized_OtherNoiseLevels_Grid_VarySize_LargeMotorNoise.py_180_2_2345_N40000_STEEPPERIODIC_UNIFORM.txt"
python3 evaluateCrossValidationResults_Synthetic_LargeMotorNoise.py "SimulateSynthetic_Parameterized_OtherNoiseLevels_Grid_VarySize_LargeMotorNoise.py_180_2_2345_N40000_UNIFORM_STEEPPERIODIC.txt"
python3 evaluateCrossValidationResults_Synthetic_LargeMotorNoise.py "SimulateSynthetic_Parameterized_OtherNoiseLevels_Grid_VarySize_LargeMotorNoise.py_180_2_2345_N40000_STEEPSHIFTED_STEEPPERIODIC.txt"
python3 evaluateCrossValidationResults_Synthetic_LargeMotorNoise.py "SimulateSynthetic_Parameterized_OtherNoiseLevels_Grid_VarySize_LargeMotorNoise.py_180_8_2345_N40000_STEEPPERIODIC_STEEPPERIODIC.txt"
python3 evaluateCrossValidationResults_Synthetic_LargeMotorNoise.py "SimulateSynthetic_Parameterized_OtherNoiseLevels_Grid_VarySize_LargeMotorNoise.py_180_8_2345_N40000_STEEPPERIODIC_UNIFORM.txt"
python3 evaluateCrossValidationResults_Synthetic_LargeMotorNoise.py "SimulateSynthetic_Parameterized_OtherNoiseLevels_Grid_VarySize_LargeMotorNoise.py_180_8_2345_N40000_UNIFORM_STEEPPERIODIC.txt"
python3 evaluateCrossValidationResults_Synthetic_LargeMotorNoise.py "SimulateSynthetic_Parameterized_OtherNoiseLevels_Grid_VarySize_LargeMotorNoise.py_180_8_2345_N40000_STEEPSHIFTED_STEEPPERIODIC.txt"


\end{comment}


\input{si-alt-simulations-medium-large-motor-noise-preview.tex}

\clearpage
\newpage

\subsection{Robustness to Regularization}\label{sec:robustness-reg}
As described in Materials and Methods, the fitting procedure from \citet{hahn2024unifying} uses regularization to encourage smooth fits for prior and encoding.
While the regularization strength can in principle be determined for each dataset individually via cross-validation, we (for computational efficiency) use a single regularization strength across all simulations, using the strength identified for the fit to the data of \citet{deGardelle2010AnOI} in \citet{hahn2024unifying}, as described in Materials and Methods.
Here, we test the robustness of identifiability to the choice of the regularization strength parameter, by fitting to the same simulated datasets as in the main simulations (four combinations of prior and encoding as in Section~\ref{sec:further-circular}, 10K trials, four noise levels), but using (i) no regularization, (ii) much weaker regularization, (iii) much stronger regularization in model fitting.
Results are shown in Figures~\ref{fig:si-alt-regularization-steepperiodic-uniform}--\ref{fig:si-alt-regularization-steepshifted-steepperiodic}.
Across datasets, prior and encoding are biased towards uniform when regularization is very strong, and become jagged when regularization is weak or absent.
Importantly, the NLL profile is very similar across regularization strengths, and the true model is identifiable even with no regularization, albeit with a more jagged fit.
This shows that regularization is useful for mitigating overfitting to a finite sample, which can lead to excessively jagged fits, but is not necessary for identifiability in principle, and that the choice of regularization strength does not qualitatively affect identifiability results.


\newcommand{\regsweeprow}[2]{%
\texttt{#2} &
\safegraphic{0.17\textwidth}{../code/Synthetic/figures/RunSynthetic_FreePrior_CosineLoss_OnSim_VIZ_OnlyModel_OtherNoiseLevels_WithGroundTruthPrior.py_SimulateSynthetic_Parameterized_OtherNoiseLevels_Grid_VarySize.py_180_2_2345_N10000_#1.txt_2_0_#2_180.pdf}{3.0cm}{fit 2 reg #2} &
\safegraphic{0.10\textwidth}{../code/Synthetic/figures/regularization_sweep_nll/SimulateSynthetic_Parameterized_OtherNoiseLevels_Grid_VarySize.py_180_2_2345_N10000_#1.txt_reg_#2_RelativeLF.pdf}{3.0cm}{nll 2 reg #2} &
\safegraphic{0.17\textwidth}{../code/Synthetic/figures/RunSynthetic_FreePrior_CosineLoss_OnSim_VIZ_OnlyModel_OtherNoiseLevels_WithGroundTruthPrior.py_SimulateSynthetic_Parameterized_OtherNoiseLevels_Grid_VarySize.py_180_8_2345_N10000_#1.txt_8_0_#2_180.pdf}{3.0cm}{fit 8 reg #2} &
\safegraphic{0.10\textwidth}{../code/Synthetic/figures/regularization_sweep_nll/SimulateSynthetic_Parameterized_OtherNoiseLevels_Grid_VarySize.py_180_8_2345_N10000_#1.txt_reg_#2_RelativeLF.pdf}{3.0cm}{nll 8 reg #2} \\
}

\newcommand{\regsweeprowlabel}[3]{%
\raisebox{0pt}[0pt][0pt]{\makebox[0.08\textwidth][c]{\textit{#3}}} &
\safegraphic{0.17\textwidth}{../code/Synthetic/figures/RunSynthetic_FreePrior_CosineLoss_OnSim_VIZ_OnlyModel_OtherNoiseLevels_WithGroundTruthPrior.py_SimulateSynthetic_Parameterized_OtherNoiseLevels_Grid_VarySize.py_180_2_2345_N10000_#1.txt_2_0_#2_180.pdf}{3.0cm}{fit 2 reg #2} &
\safegraphic{0.10\textwidth}{../code/Synthetic/figures/regularization_sweep_nll/SimulateSynthetic_Parameterized_OtherNoiseLevels_Grid_VarySize.py_180_2_2345_N10000_#1.txt_reg_#2_RelativeLF.pdf}{3.0cm}{nll 2 reg #2} &
\safegraphic{0.17\textwidth}{../code/Synthetic/figures/RunSynthetic_FreePrior_CosineLoss_OnSim_VIZ_OnlyModel_OtherNoiseLevels_WithGroundTruthPrior.py_SimulateSynthetic_Parameterized_OtherNoiseLevels_Grid_VarySize.py_180_8_2345_N10000_#1.txt_8_0_#2_180.pdf}{3.0cm}{fit 8 reg #2} &
\safegraphic{0.10\textwidth}{../code/Synthetic/figures/regularization_sweep_nll/SimulateSynthetic_Parameterized_OtherNoiseLevels_Grid_VarySize.py_180_8_2345_N10000_#1.txt_reg_#2_RelativeLF.pdf}{3.0cm}{nll 8 reg #2} \\
}

\newcommand{\regsweeptable}[1]{%
\begin{tabular}{>{\centering\arraybackslash}m{0.08\textwidth}cc@{\hspace{0.05\textwidth}}cc}
& \multicolumn{2}{c}{\textbf{Ground truth: \(p=2\)}} & \multicolumn{2}{c}{\textbf{Ground truth: \(p=8\)}} \\
\textbf{Strength} & \textbf{Model fit} & \textbf{\(\Delta\)NLL} & \textbf{Model fit} & \textbf{\(\Delta\)NLL} \\
\regsweeprowlabel{#1}{0.0}{none}
\regsweeprowlabel{#1}{1.0}{small}
\regsweeprowlabel{#1}{10.0}{standard}
\regsweeprowlabel{#1}{100.0}{large}
\end{tabular}
}

\begin{figure}[p]
\centering

\textbf{Periodic Encoding, Periodic Prior}

\ 

\regsweeptable{STEEPPERIODIC_STEEPPERIODIC}


\ 

\ 

\ 

\textbf{Uniform Encoding, Perioric Prior}

\ 

\begin{tabular}{>{\centering\arraybackslash}m{0.08\textwidth}cc@{\hspace{0.05\textwidth}}cc}
& \multicolumn{2}{c}{\textbf{Ground truth: \(p=2\)}} & \multicolumn{2}{c}{\textbf{Ground truth: \(p=8\)}} \\
\textbf{Strength} & \textbf{Model fit} & \textbf{\(\Delta\)NLL} & \textbf{Model fit} & \textbf{\(\Delta\)NLL} \\
\regsweeprowlabel{STEEPPERIODIC_UNIFORM}{0.0}{none}
\regsweeprowlabel{STEEPPERIODIC_UNIFORM}{1.0}{small}
\regsweeprowlabel{STEEPPERIODIC_UNIFORM}{10.0}{standard}
\regsweeprowlabel{STEEPPERIODIC_UNIFORM}{100.0}{large}
\end{tabular}
\caption{Robustness of model identification to the regularization strength used for fitting. Here, \textit{standard} is the regularization strength used across the circular simulations (determined by \citet{hahn2024unifying} for the behavioral data from \citet{deGardelle2010AnOI}), \textit{small} is \(10\%\) of that strength, \textit{large} is \(10\times\) that strength, and \textit{none} means zero regularization.
While the strength of regularization impacts the degree of overfitting or underfitting in the recovered prior and encoding functions, the overall shapes and the NLL profile are recovered independently of the regularization strength.
}
\label{fig:si-alt-regularization-steepperiodic-uniform}
\end{figure}

\begin{figure}[p]
\centering

\textbf{Periodic Encoding, Uniform Prior}

\ 

%\small
\begin{tabular}{>{\centering\arraybackslash}m{0.08\textwidth}cc@{\hspace{0.05\textwidth}}cc}
& \multicolumn{2}{c}{\textbf{Ground truth: \(p=2\)}} & \multicolumn{2}{c}{\textbf{Ground truth: \(p=8\)}} \\
\textbf{Strength} & \textbf{Model fit} & \textbf{\(\Delta\)NLL} & \textbf{Model fit} & \textbf{\(\Delta\)NLL} \\
\regsweeprowlabel{UNIFORM_STEEPPERIODIC}{0.0}{none}
\regsweeprowlabel{UNIFORM_STEEPPERIODIC}{1.0}{small}
\regsweeprowlabel{UNIFORM_STEEPPERIODIC}{10.0}{standard}
\regsweeprowlabel{UNIFORM_STEEPPERIODIC}{100.0}{large}
\end{tabular}

\ 

\ 

\ 

\textbf{Periodic Encoding, Shifted Prior}

\ 

%\small
\begin{tabular}{>{\centering\arraybackslash}m{0.08\textwidth}cc@{\hspace{0.05\textwidth}}cc}
& \multicolumn{2}{c}{\textbf{Ground truth: \(p=2\)}} & \multicolumn{2}{c}{\textbf{Ground truth: \(p=8\)}} \\
\textbf{Strength} & \textbf{Model fit} & \textbf{\(\Delta\)NLL} & \textbf{Model fit} & \textbf{\(\Delta\)NLL} \\
\regsweeprowlabel{STEEPSHIFTED_STEEPPERIODIC}{0.0}{none}
\regsweeprowlabel{STEEPSHIFTED_STEEPPERIODIC}{1.0}{small}
\regsweeprowlabel{STEEPSHIFTED_STEEPPERIODIC}{10.0}{standard}
\regsweeprowlabel{STEEPSHIFTED_STEEPPERIODIC}{100.0}{large}
\end{tabular}
\caption{Continuation of Figure~\ref{fig:si-alt-regularization-steepperiodic-uniform}.}
\label{fig:si-alt-regularization-steepshifted-steepperiodic}
\end{figure}


\clearpage


\subsection{Beyond Gaussian/von Mises Noise}\label{sec:other-noise}
Our theoretical results and simulations model sensory noise as Gaussian or von Mises in sensory space.
Here, we discuss the possibility of identifiability under other noise distributions.
The theoretical derivations leading to Theorems 2--4 rest on calculating moments of the noise distributions, which are key to obtaining expansions of bias and variability in the noise magnitude (Lemma~\ref{lemma:expansion}), which in turn are key to recovering model components from response distributions.
These results specifically rest on simple properties of Gaussians: odd moments are zero, and even moments are proportional to powers of the variance.
These are general properties that hold much more broadly, to any fixed-shape scale family with finite moments; this includes, e.g., Laplace distributions, which have heavier tails than Gaussians but still have finite moments.
Thus, whereas the specific algebraic form of the coefficients $\mathcal{A}, \mathcal{B}, \mathcal{C}, \mathcal{D}$ in Lemma~\ref{lemma:expansion} will depend on the specific noise distribution, our proof approach might offer a general roadmap for showing analogous identifiability results for many other noise distributions.


To illustrate this, we considered a Laplace-like sensory noise distribution in a circular sensory space, defined as
\begin{equation}\label{eq:laplace-noise}
p(m\mid \theta) \propto \exp\left(-\frac{d_{circ}(F(\theta),m)}{\sigma}\right),
\end{equation}
where $d_{circ}$ is the wrapped unsigned circular distance.

We simulated sixteen datasets (the same four prior-encoding combinations as in Section~\ref{sec:further-circular}, at $p=2$ and $p=8$, 4 noise levels, 10K or 40K trials) with this sensory noise distribution.
We show the resulting noise distributions, in comparison to those used in our main simulations, in Figure~\ref{fig:si-alt-von-mises-on-laplace-densities}.

We then applied our fitting procedure to these datasets.
In one case, we used our standard fitting procedure assuming von Mises noise; in the other case, we adapted our fitting procedure to asume Laplace-like noise (\ref{eq:laplace-noise}).
Results are shown in Figures~\ref{fig:si-alt-von-mises-on-laplace} and \ref{fig:si-alt-von-mises-on-laplace}.
In both cases, the shape of the encoding and prior are recovered appropriately in most cases; the true $L_p$ exponent has minimal or close-to-minimal NLL.
In the misspecified case (fitting assuming von Mises noise), there are a few cases of incompletely identified loss function (e.g., at prioridic encoding and shifted prior at $L_2$ loss, NLL is low for $p=2,4,6,8$) or mismatch in the shape of the prior (periodic encoding and uniform prior at $L_8$ loss; the recovered prior appears to have bumps).
In the correctly specified case (fitting asuming Laplace-like noise), prior and encoding are recovered closely throughout, at least at 40K trials.
Taken together, these results suggest that our model fitting approach can be successful beyond Gaussian/von Mises noise on data from other symmetric noise distributions, and that even fitting under misspecification of the noise distribution may have rather limited effect on model identifiability.

Beyond symmetric noise, there may be situations where noise could be skewed or otherwise asymmetric in sensory space.
Such noise distributions might lead to further challenges, because the skewness may enter the bias, which could make the skewness potentially entangled with the prior.
Formally, for such distributions, the odd moments may be nonzero; while the overall approach based on Lemma~\ref{lemma:expansion} might still be viable, the increased complexity might lead to additional challenges in analysis. We leave understanding of skewed noise distributions to future work.

\begin{figure}[ht]
\centering
\safegraphic{0.49\textwidth}{figures/von_mises_on_laplace_density_von_mises.pdf}{4.0cm}{von mises densities}
\hfill
\safegraphic{0.49\textwidth}{figures/von_mises_on_laplace_density_laplace.pdf}{4.0cm}{laplace densities}
\caption{Sensory-noise densities under the von Mises family (panel A) and the Laplace-like family in Equation~\ref{eq:laplace-noise} (panel B), shown for the four noise levels used in our circular simulations. The code now follows the same orientation convention as the simulation figures: the internal noise variable lives on the full circle, but plotted orientation errors are the corresponding circular differences divided by \(2\).}
\label{fig:si-alt-von-mises-on-laplace-densities}
\end{figure}



\begin{figure}[ht]
\centering

{\textbf{10,000 Trials}}

\begin{tabular}{>{\centering\arraybackslash}m{0.17\textwidth}cccc}
& \multicolumn{2}{c}{\textbf{Ground truth: \(p=2\)}} & \multicolumn{2}{c}{\textbf{Ground truth: \(p=8\)}} \\
& Fit at \(p=2\) & \(\Delta\)NLL & Fit at \(p=8\) & \(\Delta\)NLL \\
\safegraphic{0.14\textwidth}{../code/Synthetic/figures/CounterfactualModel_VIZ_ByNoiseMagnGauge_Components.py_2345_STEEPPERIODIC_STEEPPERIODIC_2_0_10.0_180.pdf}{2.2cm}{gt periodic periodic} &
\safegraphic{0.18\textwidth}{../code/Synthetic/figures/RunSynthetic_FreePrior_CosineLoss_OnSim_VIZ_OnlyModel_OtherNoiseLevels.py_SimulateSynthetic_Parameterized_OtherNoiseLevels_Grid_VarySize_LaplaceNoise.py_180_2_2345_N10000_STEEPPERIODIC_STEEPPERIODIC.txt_2_0_10.0_180.pdf}{3.0cm}{vm fit p=2} &
\safegraphic{0.10\textwidth}{../code/Synthetic/figures/evaluateCrossValidationResults_Synthetic_Gardelle_NonF.py_SimulateSynthetic_Parameterized_OtherNoiseLevels_Grid_VarySize_LaplaceNoise.py_180_2_2345_N10000_STEEPPERIODIC_STEEPPERIODIC.txt_RelativeLF.pdf}{3.0cm}{NLL p=2} &
\safegraphic{0.18\textwidth}{../code/Synthetic/figures/RunSynthetic_FreePrior_CosineLoss_OnSim_VIZ_OnlyModel_OtherNoiseLevels.py_SimulateSynthetic_Parameterized_OtherNoiseLevels_Grid_VarySize_LaplaceNoise.py_180_8_2345_N10000_STEEPPERIODIC_STEEPPERIODIC.txt_8_0_10.0_180.pdf}{3.0cm}{vm fit p=8} &
\safegraphic{0.10\textwidth}{../code/Synthetic/figures/evaluateCrossValidationResults_Synthetic_Gardelle_NonF.py_SimulateSynthetic_Parameterized_OtherNoiseLevels_Grid_VarySize_LaplaceNoise.py_180_8_2345_N10000_STEEPPERIODIC_STEEPPERIODIC.txt_RelativeLF.pdf}{3.0cm}{NLL p=8} \\
\safegraphic{0.14\textwidth}{../code/Synthetic/figures/CounterfactualModel_VIZ_ByNoiseMagnGauge_Components.py_2345_STEEPPERIODIC_UNIFORM_2_0_10.0_180.pdf}{2.2cm}{gt periodic uniform} &
\safegraphic{0.18\textwidth}{../code/Synthetic/figures/RunSynthetic_FreePrior_CosineLoss_OnSim_VIZ_OnlyModel_OtherNoiseLevels.py_SimulateSynthetic_Parameterized_OtherNoiseLevels_Grid_VarySize_LaplaceNoise.py_180_2_2345_N10000_STEEPPERIODIC_UNIFORM.txt_2_0_10.0_180.pdf}{3.0cm}{vm fit p=2} &
\safegraphic{0.10\textwidth}{../code/Synthetic/figures/evaluateCrossValidationResults_Synthetic_Gardelle_NonF.py_SimulateSynthetic_Parameterized_OtherNoiseLevels_Grid_VarySize_LaplaceNoise.py_180_2_2345_N10000_STEEPPERIODIC_UNIFORM.txt_RelativeLF.pdf}{3.0cm}{NLL p=2} &
\safegraphic{0.18\textwidth}{../code/Synthetic/figures/RunSynthetic_FreePrior_CosineLoss_OnSim_VIZ_OnlyModel_OtherNoiseLevels.py_SimulateSynthetic_Parameterized_OtherNoiseLevels_Grid_VarySize_LaplaceNoise.py_180_8_2345_N10000_STEEPPERIODIC_UNIFORM.txt_8_0_10.0_180.pdf}{3.0cm}{vm fit p=8} &
\safegraphic{0.10\textwidth}{../code/Synthetic/figures/evaluateCrossValidationResults_Synthetic_Gardelle_NonF.py_SimulateSynthetic_Parameterized_OtherNoiseLevels_Grid_VarySize_LaplaceNoise.py_180_8_2345_N10000_STEEPPERIODIC_UNIFORM.txt_RelativeLF.pdf}{3.0cm}{NLL p=8} \\
\safegraphic{0.14\textwidth}{../code/Synthetic/figures/CounterfactualModel_VIZ_ByNoiseMagnGauge_Components.py_2345_UNIFORM_STEEPPERIODIC_2_0_10.0_180.pdf}{2.2cm}{gt uniform periodic} &
\safegraphic{0.18\textwidth}{../code/Synthetic/figures/RunSynthetic_FreePrior_CosineLoss_OnSim_VIZ_OnlyModel_OtherNoiseLevels.py_SimulateSynthetic_Parameterized_OtherNoiseLevels_Grid_VarySize_LaplaceNoise.py_180_2_2345_N10000_UNIFORM_STEEPPERIODIC.txt_2_0_10.0_180.pdf}{3.0cm}{vm fit p=2} &
\safegraphic{0.10\textwidth}{../code/Synthetic/figures/evaluateCrossValidationResults_Synthetic_Gardelle_NonF.py_SimulateSynthetic_Parameterized_OtherNoiseLevels_Grid_VarySize_LaplaceNoise.py_180_2_2345_N10000_UNIFORM_STEEPPERIODIC.txt_RelativeLF.pdf}{3.0cm}{NLL p=2} &
\safegraphic{0.18\textwidth}{../code/Synthetic/figures/RunSynthetic_FreePrior_CosineLoss_OnSim_VIZ_OnlyModel_OtherNoiseLevels.py_SimulateSynthetic_Parameterized_OtherNoiseLevels_Grid_VarySize_LaplaceNoise.py_180_8_2345_N10000_UNIFORM_STEEPPERIODIC.txt_8_0_10.0_180.pdf}{3.0cm}{vm fit p=8} &
\safegraphic{0.10\textwidth}{../code/Synthetic/figures/evaluateCrossValidationResults_Synthetic_Gardelle_NonF.py_SimulateSynthetic_Parameterized_OtherNoiseLevels_Grid_VarySize_LaplaceNoise.py_180_8_2345_N10000_UNIFORM_STEEPPERIODIC.txt_RelativeLF.pdf}{3.0cm}{NLL p=8} \\
\safegraphic{0.14\textwidth}{../code/Synthetic/figures/CounterfactualModel_VIZ_ByNoiseMagnGauge_Components.py_2345_STEEPSHIFTED_STEEPPERIODIC_2_0_10.0_180.pdf}{2.2cm}{gt shifted periodic} &
\safegraphic{0.18\textwidth}{../code/Synthetic/figures/RunSynthetic_FreePrior_CosineLoss_OnSim_VIZ_OnlyModel_OtherNoiseLevels.py_SimulateSynthetic_Parameterized_OtherNoiseLevels_Grid_VarySize_LaplaceNoise.py_180_2_2345_N10000_STEEPSHIFTED_STEEPPERIODIC.txt_2_0_10.0_180.pdf}{3.0cm}{vm fit p=2} &
\safegraphic{0.10\textwidth}{../code/Synthetic/figures/evaluateCrossValidationResults_Synthetic_Gardelle_NonF.py_SimulateSynthetic_Parameterized_OtherNoiseLevels_Grid_VarySize_LaplaceNoise.py_180_2_2345_N10000_STEEPSHIFTED_STEEPPERIODIC.txt_RelativeLF.pdf}{3.0cm}{NLL p=2} &
\safegraphic{0.18\textwidth}{../code/Synthetic/figures/RunSynthetic_FreePrior_CosineLoss_OnSim_VIZ_OnlyModel_OtherNoiseLevels.py_SimulateSynthetic_Parameterized_OtherNoiseLevels_Grid_VarySize_LaplaceNoise.py_180_8_2345_N10000_STEEPSHIFTED_STEEPPERIODIC.txt_8_0_10.0_180.pdf}{3.0cm}{vm fit p=8} &
\safegraphic{0.10\textwidth}{../code/Synthetic/figures/evaluateCrossValidationResults_Synthetic_Gardelle_NonF.py_SimulateSynthetic_Parameterized_OtherNoiseLevels_Grid_VarySize_LaplaceNoise.py_180_8_2345_N10000_STEEPSHIFTED_STEEPPERIODIC.txt_RelativeLF.pdf}{3.0cm}{NLL p=8}
\end{tabular}

\ 

\ 

\ 

{\textbf{40,000 Trials}}

\begin{tabular}{>{\centering\arraybackslash}m{0.17\textwidth}cccc}
& \multicolumn{2}{c}{\textbf{Ground truth: \(p=2\)}} & \multicolumn{2}{c}{\textbf{Ground truth: \(p=8\)}} \\
& Fit at \(p=2\) & \(\Delta\)NLL & Fit at \(p=8\) & \(\Delta\)NLL \\
\safegraphic{0.14\textwidth}{../code/Synthetic/figures/CounterfactualModel_VIZ_ByNoiseMagnGauge_Components.py_2345_STEEPPERIODIC_STEEPPERIODIC_2_0_10.0_180.pdf}{2.2cm}{gt periodic periodic} &
\safegraphic{0.18\textwidth}{../code/Synthetic/figures/RunSynthetic_FreePrior_CosineLoss_OnSim_VIZ_OnlyModel_OtherNoiseLevels.py_SimulateSynthetic_Parameterized_OtherNoiseLevels_Grid_VarySize_LaplaceNoise.py_180_2_2345_N40000_STEEPPERIODIC_STEEPPERIODIC.txt_2_0_10.0_180.pdf}{3.0cm}{vm fit p=2 40k} &
\safegraphic{0.10\textwidth}{../code/Synthetic/figures/evaluateCrossValidationResults_Synthetic_Gardelle_NonF.py_SimulateSynthetic_Parameterized_OtherNoiseLevels_Grid_VarySize_LaplaceNoise.py_180_2_2345_N40000_STEEPPERIODIC_STEEPPERIODIC.txt_RelativeLF.pdf}{3.0cm}{NLL p=2 40k} &
\safegraphic{0.18\textwidth}{../code/Synthetic/figures/RunSynthetic_FreePrior_CosineLoss_OnSim_VIZ_OnlyModel_OtherNoiseLevels.py_SimulateSynthetic_Parameterized_OtherNoiseLevels_Grid_VarySize_LaplaceNoise.py_180_8_2345_N40000_STEEPPERIODIC_STEEPPERIODIC.txt_8_0_10.0_180.pdf}{3.0cm}{vm fit p=8 40k} &
\safegraphic{0.10\textwidth}{../code/Synthetic/figures/evaluateCrossValidationResults_Synthetic_Gardelle_NonF.py_SimulateSynthetic_Parameterized_OtherNoiseLevels_Grid_VarySize_LaplaceNoise.py_180_8_2345_N40000_STEEPPERIODIC_STEEPPERIODIC.txt_RelativeLF.pdf}{3.0cm}{NLL p=8 40k} \\
\safegraphic{0.14\textwidth}{../code/Synthetic/figures/CounterfactualModel_VIZ_ByNoiseMagnGauge_Components.py_2345_STEEPPERIODIC_UNIFORM_2_0_10.0_180.pdf}{2.2cm}{gt periodic uniform} &
\safegraphic{0.18\textwidth}{../code/Synthetic/figures/RunSynthetic_FreePrior_CosineLoss_OnSim_VIZ_OnlyModel_OtherNoiseLevels.py_SimulateSynthetic_Parameterized_OtherNoiseLevels_Grid_VarySize_LaplaceNoise.py_180_2_2345_N40000_STEEPPERIODIC_UNIFORM.txt_2_0_10.0_180.pdf}{3.0cm}{vm fit p=2 40k} &
\safegraphic{0.10\textwidth}{../code/Synthetic/figures/evaluateCrossValidationResults_Synthetic_Gardelle_NonF.py_SimulateSynthetic_Parameterized_OtherNoiseLevels_Grid_VarySize_LaplaceNoise.py_180_2_2345_N40000_STEEPPERIODIC_UNIFORM.txt_RelativeLF.pdf}{3.0cm}{NLL p=2 40k} &
\safegraphic{0.18\textwidth}{../code/Synthetic/figures/RunSynthetic_FreePrior_CosineLoss_OnSim_VIZ_OnlyModel_OtherNoiseLevels.py_SimulateSynthetic_Parameterized_OtherNoiseLevels_Grid_VarySize_LaplaceNoise.py_180_8_2345_N40000_STEEPPERIODIC_UNIFORM.txt_8_0_10.0_180.pdf}{3.0cm}{vm fit p=8 40k} &
\safegraphic{0.10\textwidth}{../code/Synthetic/figures/evaluateCrossValidationResults_Synthetic_Gardelle_NonF.py_SimulateSynthetic_Parameterized_OtherNoiseLevels_Grid_VarySize_LaplaceNoise.py_180_8_2345_N40000_STEEPPERIODIC_UNIFORM.txt_RelativeLF.pdf}{3.0cm}{NLL p=8 40k} \\
\safegraphic{0.14\textwidth}{../code/Synthetic/figures/CounterfactualModel_VIZ_ByNoiseMagnGauge_Components.py_2345_UNIFORM_STEEPPERIODIC_2_0_10.0_180.pdf}{2.2cm}{gt uniform periodic} &
\safegraphic{0.18\textwidth}{../code/Synthetic/figures/RunSynthetic_FreePrior_CosineLoss_OnSim_VIZ_OnlyModel_OtherNoiseLevels.py_SimulateSynthetic_Parameterized_OtherNoiseLevels_Grid_VarySize_LaplaceNoise.py_180_2_2345_N40000_UNIFORM_STEEPPERIODIC.txt_2_0_10.0_180.pdf}{3.0cm}{vm fit p=2 40k} &
\safegraphic{0.10\textwidth}{../code/Synthetic/figures/evaluateCrossValidationResults_Synthetic_Gardelle_NonF.py_SimulateSynthetic_Parameterized_OtherNoiseLevels_Grid_VarySize_LaplaceNoise.py_180_2_2345_N40000_UNIFORM_STEEPPERIODIC.txt_RelativeLF.pdf}{3.0cm}{NLL p=2 40k} &
\safegraphic{0.18\textwidth}{../code/Synthetic/figures/RunSynthetic_FreePrior_CosineLoss_OnSim_VIZ_OnlyModel_OtherNoiseLevels.py_SimulateSynthetic_Parameterized_OtherNoiseLevels_Grid_VarySize_LaplaceNoise.py_180_8_2345_N40000_UNIFORM_STEEPPERIODIC.txt_8_0_10.0_180.pdf}{3.0cm}{vm fit p=8 40k} &
\safegraphic{0.10\textwidth}{../code/Synthetic/figures/evaluateCrossValidationResults_Synthetic_Gardelle_NonF.py_SimulateSynthetic_Parameterized_OtherNoiseLevels_Grid_VarySize_LaplaceNoise.py_180_8_2345_N40000_UNIFORM_STEEPPERIODIC.txt_RelativeLF.pdf}{3.0cm}{NLL p=8 40k} \\
\safegraphic{0.14\textwidth}{../code/Synthetic/figures/CounterfactualModel_VIZ_ByNoiseMagnGauge_Components.py_2345_STEEPSHIFTED_STEEPPERIODIC_2_0_10.0_180.pdf}{2.2cm}{gt shifted periodic} &
\safegraphic{0.18\textwidth}{../code/Synthetic/figures/RunSynthetic_FreePrior_CosineLoss_OnSim_VIZ_OnlyModel_OtherNoiseLevels.py_SimulateSynthetic_Parameterized_OtherNoiseLevels_Grid_VarySize_LaplaceNoise.py_180_2_2345_N40000_STEEPSHIFTED_STEEPPERIODIC.txt_2_0_10.0_180.pdf}{3.0cm}{vm fit p=2 40k} &
\safegraphic{0.10\textwidth}{../code/Synthetic/figures/evaluateCrossValidationResults_Synthetic_Gardelle_NonF.py_SimulateSynthetic_Parameterized_OtherNoiseLevels_Grid_VarySize_LaplaceNoise.py_180_2_2345_N40000_STEEPSHIFTED_STEEPPERIODIC.txt_RelativeLF.pdf}{3.0cm}{NLL p=2 40k} &
\safegraphic{0.18\textwidth}{../code/Synthetic/figures/RunSynthetic_FreePrior_CosineLoss_OnSim_VIZ_OnlyModel_OtherNoiseLevels.py_SimulateSynthetic_Parameterized_OtherNoiseLevels_Grid_VarySize_LaplaceNoise.py_180_8_2345_N40000_STEEPSHIFTED_STEEPPERIODIC.txt_8_0_10.0_180.pdf}{3.0cm}{vm fit p=8 40k} &
\safegraphic{0.10\textwidth}{../code/Synthetic/figures/evaluateCrossValidationResults_Synthetic_Gardelle_NonF.py_SimulateSynthetic_Parameterized_OtherNoiseLevels_Grid_VarySize_LaplaceNoise.py_180_8_2345_N40000_STEEPSHIFTED_STEEPPERIODIC.txt_RelativeLF.pdf}{3.0cm}{NLL p=8 40k}
\end{tabular}
\caption{Applying our fitting procedure (assuming von Mises noise) to data generated under Laplace-like sensory noise.}
\label{fig:si-alt-von-mises-on-laplace}
\end{figure}



\begin{comment}
python3 RunSynthetic_FreePrior_CosineLoss_OnSim_VIZ_OnlyModel_OtherNoiseLevels.py 2 0 10.0 180 "SimulateSynthetic_Parameterized_OtherNoiseLevels_Grid_VarySize_LaplaceNoise.py_180_2_2345_N10000_STEEPPERIODIC_STEEPPERIODIC.txt"
python3 RunSynthetic_FreePrior_CosineLoss_OnSim_VIZ_OnlyModel_OtherNoiseLevels.py 2 0 10.0 180 "SimulateSynthetic_Parameterized_OtherNoiseLevels_Grid_VarySize_LaplaceNoise.py_180_2_2345_N10000_STEEPPERIODIC_UNIFORM.txt"
python3 RunSynthetic_FreePrior_CosineLoss_OnSim_VIZ_OnlyModel_OtherNoiseLevels.py 2 0 10.0 180 "SimulateSynthetic_Parameterized_OtherNoiseLevels_Grid_VarySize_LaplaceNoise.py_180_2_2345_N10000_UNIFORM_STEEPPERIODIC.txt"
python3 RunSynthetic_FreePrior_CosineLoss_OnSim_VIZ_OnlyModel_OtherNoiseLevels.py 2 0 10.0 180 "SimulateSynthetic_Parameterized_OtherNoiseLevels_Grid_VarySize_LaplaceNoise.py_180_2_2345_N10000_STEEPSHIFTED_STEEPPERIODIC.txt"
python3 RunSynthetic_FreePrior_CosineLoss_OnSim_VIZ_OnlyModel_OtherNoiseLevels.py 8 0 10.0 180 "SimulateSynthetic_Parameterized_OtherNoiseLevels_Grid_VarySize_LaplaceNoise.py_180_8_2345_N10000_STEEPPERIODIC_STEEPPERIODIC.txt"
python3 RunSynthetic_FreePrior_CosineLoss_OnSim_VIZ_OnlyModel_OtherNoiseLevels.py 8 0 10.0 180 "SimulateSynthetic_Parameterized_OtherNoiseLevels_Grid_VarySize_LaplaceNoise.py_180_8_2345_N10000_STEEPPERIODIC_UNIFORM.txt"
python3 RunSynthetic_FreePrior_CosineLoss_OnSim_VIZ_OnlyModel_OtherNoiseLevels.py 8 0 10.0 180 "SimulateSynthetic_Parameterized_OtherNoiseLevels_Grid_VarySize_LaplaceNoise.py_180_8_2345_N10000_UNIFORM_STEEPPERIODIC.txt"
python3 RunSynthetic_FreePrior_CosineLoss_OnSim_VIZ_OnlyModel_OtherNoiseLevels.py 8 0 10.0 180 "SimulateSynthetic_Parameterized_OtherNoiseLevels_Grid_VarySize_LaplaceNoise.py_180_8_2345_N10000_STEEPSHIFTED_STEEPPERIODIC.txt"

python3 evaluateCrossValidationResults_Synthetic_Gardelle_NonF.py "SimulateSynthetic_Parameterized_OtherNoiseLevels_Grid_VarySize_LaplaceNoise.py_180_2_2345_N10000_STEEPPERIODIC_STEEPPERIODIC.txt"
python3 evaluateCrossValidationResults_Synthetic_Gardelle_NonF.py "SimulateSynthetic_Parameterized_OtherNoiseLevels_Grid_VarySize_LaplaceNoise.py_180_2_2345_N10000_STEEPPERIODIC_UNIFORM.txt"
python3 evaluateCrossValidationResults_Synthetic_Gardelle_NonF.py "SimulateSynthetic_Parameterized_OtherNoiseLevels_Grid_VarySize_LaplaceNoise.py_180_2_2345_N10000_UNIFORM_STEEPPERIODIC.txt"
python3 evaluateCrossValidationResults_Synthetic_Gardelle_NonF.py "SimulateSynthetic_Parameterized_OtherNoiseLevels_Grid_VarySize_LaplaceNoise.py_180_2_2345_N10000_STEEPSHIFTED_STEEPPERIODIC.txt"
python3 evaluateCrossValidationResults_Synthetic_Gardelle_NonF.py "SimulateSynthetic_Parameterized_OtherNoiseLevels_Grid_VarySize_LaplaceNoise.py_180_8_2345_N10000_STEEPPERIODIC_STEEPPERIODIC.txt"
python3 evaluateCrossValidationResults_Synthetic_Gardelle_NonF.py "SimulateSynthetic_Parameterized_OtherNoiseLevels_Grid_VarySize_LaplaceNoise.py_180_8_2345_N10000_STEEPPERIODIC_UNIFORM.txt"
python3 evaluateCrossValidationResults_Synthetic_Gardelle_NonF.py "SimulateSynthetic_Parameterized_OtherNoiseLevels_Grid_VarySize_LaplaceNoise.py_180_8_2345_N10000_UNIFORM_STEEPPERIODIC.txt"
python3 evaluateCrossValidationResults_Synthetic_Gardelle_NonF.py "SimulateSynthetic_Parameterized_OtherNoiseLevels_Grid_VarySize_LaplaceNoise.py_180_8_2345_N10000_STEEPSHIFTED_STEEPPERIODIC.txt"

cd ../code/Synthetic
bash run_viz_von_mises_on_laplace_noise_40000.sh
bash run_evaluate_von_mises_on_laplace_noise_40000.sh

bash run_training_laplace_noise_40000.sh
bash run_viz_laplace_noise_40000.sh
bash run_evaluate_laplace_noise_40000.sh
\end{comment}


\begin{figure}[ht]
\centering

{\textbf{10,000 Trials}}

\begin{tabular}{>{\raggedright\arraybackslash}m{0.17\textwidth}cccc}
& \multicolumn{2}{c}{\textbf{Ground truth: \(p=2\)}} & \multicolumn{2}{c}{\textbf{Ground truth: \(p=8\)}} \\
& Fit at \(p=2\) & \(\Delta\)NLL & Fit at \(p=8\) & \(\Delta\)NLL \\
\safegraphic{0.14\textwidth}{../code/Synthetic/figures/CounterfactualModel_VIZ_ByNoiseMagnGauge_Components.py_2345_STEEPPERIODIC_STEEPPERIODIC_2_0_10.0_180.pdf}{2.2cm}{gt periodic periodic} &
\safegraphic{0.18\textwidth}{../code/Synthetic/figures/RunSynthetic_FreePrior_CosineLoss_OnSim_LaplaceNoise_VIZ_OnlyModel_OtherNoiseLevels.py_SimulateSynthetic_Parameterized_OtherNoiseLevels_Grid_VarySize_LaplaceNoise.py_180_2_2345_N10000_STEEPPERIODIC_STEEPPERIODIC.txt_2_0_10.0_180.pdf}{3.0cm}{laplace fit p=2} &
\safegraphic{0.10\textwidth}{../code/Synthetic/figures/evaluateCrossValidationResults_Synthetic_LaplaceNoise.py_SimulateSynthetic_Parameterized_OtherNoiseLevels_Grid_VarySize_LaplaceNoise.py_180_2_2345_N10000_STEEPPERIODIC_STEEPPERIODIC.txt_RelativeLF.pdf}{3.0cm}{NLL p=2} &
\safegraphic{0.18\textwidth}{../code/Synthetic/figures/RunSynthetic_FreePrior_CosineLoss_OnSim_LaplaceNoise_VIZ_OnlyModel_OtherNoiseLevels.py_SimulateSynthetic_Parameterized_OtherNoiseLevels_Grid_VarySize_LaplaceNoise.py_180_8_2345_N10000_STEEPPERIODIC_STEEPPERIODIC.txt_8_0_10.0_180.pdf}{3.0cm}{laplace fit p=8} &
\safegraphic{0.10\textwidth}{../code/Synthetic/figures/evaluateCrossValidationResults_Synthetic_LaplaceNoise.py_SimulateSynthetic_Parameterized_OtherNoiseLevels_Grid_VarySize_LaplaceNoise.py_180_8_2345_N10000_STEEPPERIODIC_STEEPPERIODIC.txt_RelativeLF.pdf}{3.0cm}{NLL p=8} \\
\safegraphic{0.14\textwidth}{../code/Synthetic/figures/CounterfactualModel_VIZ_ByNoiseMagnGauge_Components.py_2345_STEEPPERIODIC_UNIFORM_2_0_10.0_180.pdf}{2.2cm}{gt periodic uniform} &
\safegraphic{0.18\textwidth}{../code/Synthetic/figures/RunSynthetic_FreePrior_CosineLoss_OnSim_LaplaceNoise_VIZ_OnlyModel_OtherNoiseLevels.py_SimulateSynthetic_Parameterized_OtherNoiseLevels_Grid_VarySize_LaplaceNoise.py_180_2_2345_N10000_STEEPPERIODIC_UNIFORM.txt_2_0_10.0_180.pdf}{3.0cm}{laplace fit p=2} &
\safegraphic{0.10\textwidth}{../code/Synthetic/figures/evaluateCrossValidationResults_Synthetic_LaplaceNoise.py_SimulateSynthetic_Parameterized_OtherNoiseLevels_Grid_VarySize_LaplaceNoise.py_180_2_2345_N10000_STEEPPERIODIC_UNIFORM.txt_RelativeLF.pdf}{3.0cm}{NLL p=2} &
\safegraphic{0.18\textwidth}{../code/Synthetic/figures/RunSynthetic_FreePrior_CosineLoss_OnSim_LaplaceNoise_VIZ_OnlyModel_OtherNoiseLevels.py_SimulateSynthetic_Parameterized_OtherNoiseLevels_Grid_VarySize_LaplaceNoise.py_180_8_2345_N10000_STEEPPERIODIC_UNIFORM.txt_8_0_10.0_180.pdf}{3.0cm}{laplace fit p=8} &
\safegraphic{0.10\textwidth}{../code/Synthetic/figures/evaluateCrossValidationResults_Synthetic_LaplaceNoise.py_SimulateSynthetic_Parameterized_OtherNoiseLevels_Grid_VarySize_LaplaceNoise.py_180_8_2345_N10000_STEEPPERIODIC_UNIFORM.txt_RelativeLF.pdf}{3.0cm}{NLL p=8} \\
\safegraphic{0.14\textwidth}{../code/Synthetic/figures/CounterfactualModel_VIZ_ByNoiseMagnGauge_Components.py_2345_UNIFORM_STEEPPERIODIC_2_0_10.0_180.pdf}{2.2cm}{gt uniform periodic} &
\safegraphic{0.18\textwidth}{../code/Synthetic/figures/RunSynthetic_FreePrior_CosineLoss_OnSim_LaplaceNoise_VIZ_OnlyModel_OtherNoiseLevels.py_SimulateSynthetic_Parameterized_OtherNoiseLevels_Grid_VarySize_LaplaceNoise.py_180_2_2345_N10000_UNIFORM_STEEPPERIODIC.txt_2_0_10.0_180.pdf}{3.0cm}{laplace fit p=2} &
\safegraphic{0.10\textwidth}{../code/Synthetic/figures/evaluateCrossValidationResults_Synthetic_LaplaceNoise.py_SimulateSynthetic_Parameterized_OtherNoiseLevels_Grid_VarySize_LaplaceNoise.py_180_2_2345_N10000_UNIFORM_STEEPPERIODIC.txt_RelativeLF.pdf}{3.0cm}{NLL p=2} &
\safegraphic{0.18\textwidth}{../code/Synthetic/figures/RunSynthetic_FreePrior_CosineLoss_OnSim_LaplaceNoise_VIZ_OnlyModel_OtherNoiseLevels.py_SimulateSynthetic_Parameterized_OtherNoiseLevels_Grid_VarySize_LaplaceNoise.py_180_8_2345_N10000_UNIFORM_STEEPPERIODIC.txt_8_0_10.0_180.pdf}{3.0cm}{laplace fit p=8} &
\safegraphic{0.10\textwidth}{../code/Synthetic/figures/evaluateCrossValidationResults_Synthetic_LaplaceNoise.py_SimulateSynthetic_Parameterized_OtherNoiseLevels_Grid_VarySize_LaplaceNoise.py_180_8_2345_N10000_UNIFORM_STEEPPERIODIC.txt_RelativeLF.pdf}{3.0cm}{NLL p=8} \\
\safegraphic{0.14\textwidth}{../code/Synthetic/figures/CounterfactualModel_VIZ_ByNoiseMagnGauge_Components.py_2345_STEEPSHIFTED_STEEPPERIODIC_2_0_10.0_180.pdf}{2.2cm}{gt shifted periodic} &
\safegraphic{0.18\textwidth}{../code/Synthetic/figures/RunSynthetic_FreePrior_CosineLoss_OnSim_LaplaceNoise_VIZ_OnlyModel_OtherNoiseLevels.py_SimulateSynthetic_Parameterized_OtherNoiseLevels_Grid_VarySize_LaplaceNoise.py_180_2_2345_N10000_STEEPSHIFTED_STEEPPERIODIC.txt_2_0_10.0_180.pdf}{3.0cm}{laplace fit p=2} &
\safegraphic{0.10\textwidth}{../code/Synthetic/figures/evaluateCrossValidationResults_Synthetic_LaplaceNoise.py_SimulateSynthetic_Parameterized_OtherNoiseLevels_Grid_VarySize_LaplaceNoise.py_180_2_2345_N10000_STEEPSHIFTED_STEEPPERIODIC.txt_RelativeLF.pdf}{3.0cm}{NLL p=2} &
\safegraphic{0.18\textwidth}{../code/Synthetic/figures/RunSynthetic_FreePrior_CosineLoss_OnSim_LaplaceNoise_VIZ_OnlyModel_OtherNoiseLevels.py_SimulateSynthetic_Parameterized_OtherNoiseLevels_Grid_VarySize_LaplaceNoise.py_180_8_2345_N10000_STEEPSHIFTED_STEEPPERIODIC.txt_8_0_10.0_180.pdf}{3.0cm}{laplace fit p=8} &
\safegraphic{0.10\textwidth}{../code/Synthetic/figures/evaluateCrossValidationResults_Synthetic_LaplaceNoise.py_SimulateSynthetic_Parameterized_OtherNoiseLevels_Grid_VarySize_LaplaceNoise.py_180_8_2345_N10000_STEEPSHIFTED_STEEPPERIODIC.txt_RelativeLF.pdf}{3.0cm}{NLL p=8}
\end{tabular}

\ 

\ 

\ 

{\textbf{40,000 Trials}}

\begin{tabular}{>{\raggedright\arraybackslash}m{0.17\textwidth}cccc}
& \multicolumn{2}{c}{\textbf{Ground truth: \(p=2\)}} & \multicolumn{2}{c}{\textbf{Ground truth: \(p=8\)}} \\
& Fit at \(p=2\) & \(\Delta\)NLL & Fit at \(p=8\) & \(\Delta\)NLL \\
\safegraphic{0.14\textwidth}{../code/Synthetic/figures/CounterfactualModel_VIZ_ByNoiseMagnGauge_Components.py_2345_STEEPPERIODIC_STEEPPERIODIC_2_0_10.0_180.pdf}{2.2cm}{gt periodic periodic} &
\safegraphic{0.18\textwidth}{../code/Synthetic/figures/RunSynthetic_FreePrior_CosineLoss_OnSim_LaplaceNoise_VIZ_OnlyModel_OtherNoiseLevels.py_SimulateSynthetic_Parameterized_OtherNoiseLevels_Grid_VarySize_LaplaceNoise.py_180_2_2345_N40000_STEEPPERIODIC_STEEPPERIODIC.txt_2_0_10.0_180.pdf}{3.0cm}{laplace fit p=2 40k} &
\safegraphic{0.10\textwidth}{../code/Synthetic/figures/evaluateCrossValidationResults_Synthetic_LaplaceNoise.py_SimulateSynthetic_Parameterized_OtherNoiseLevels_Grid_VarySize_LaplaceNoise.py_180_2_2345_N40000_STEEPPERIODIC_STEEPPERIODIC.txt_RelativeLF.pdf}{3.0cm}{NLL p=2 40k} &
\safegraphic{0.18\textwidth}{../code/Synthetic/figures/RunSynthetic_FreePrior_CosineLoss_OnSim_LaplaceNoise_VIZ_OnlyModel_OtherNoiseLevels.py_SimulateSynthetic_Parameterized_OtherNoiseLevels_Grid_VarySize_LaplaceNoise.py_180_8_2345_N40000_STEEPPERIODIC_STEEPPERIODIC.txt_8_0_10.0_180.pdf}{3.0cm}{laplace fit p=8 40k} &
\safegraphic{0.10\textwidth}{../code/Synthetic/figures/evaluateCrossValidationResults_Synthetic_LaplaceNoise.py_SimulateSynthetic_Parameterized_OtherNoiseLevels_Grid_VarySize_LaplaceNoise.py_180_8_2345_N40000_STEEPPERIODIC_STEEPPERIODIC.txt_RelativeLF.pdf}{3.0cm}{NLL p=8 40k} \\
\safegraphic{0.14\textwidth}{../code/Synthetic/figures/CounterfactualModel_VIZ_ByNoiseMagnGauge_Components.py_2345_STEEPPERIODIC_UNIFORM_2_0_10.0_180.pdf}{2.2cm}{gt periodic uniform} &
\safegraphic{0.18\textwidth}{../code/Synthetic/figures/RunSynthetic_FreePrior_CosineLoss_OnSim_LaplaceNoise_VIZ_OnlyModel_OtherNoiseLevels.py_SimulateSynthetic_Parameterized_OtherNoiseLevels_Grid_VarySize_LaplaceNoise.py_180_2_2345_N40000_STEEPPERIODIC_UNIFORM.txt_2_0_10.0_180.pdf}{3.0cm}{laplace fit p=2 40k} &
\safegraphic{0.10\textwidth}{../code/Synthetic/figures/evaluateCrossValidationResults_Synthetic_LaplaceNoise.py_SimulateSynthetic_Parameterized_OtherNoiseLevels_Grid_VarySize_LaplaceNoise.py_180_2_2345_N40000_STEEPPERIODIC_UNIFORM.txt_RelativeLF.pdf}{3.0cm}{NLL p=2 40k} &
\safegraphic{0.18\textwidth}{../code/Synthetic/figures/RunSynthetic_FreePrior_CosineLoss_OnSim_LaplaceNoise_VIZ_OnlyModel_OtherNoiseLevels.py_SimulateSynthetic_Parameterized_OtherNoiseLevels_Grid_VarySize_LaplaceNoise.py_180_8_2345_N40000_STEEPPERIODIC_UNIFORM.txt_8_0_10.0_180.pdf}{3.0cm}{laplace fit p=8 40k} &
\safegraphic{0.10\textwidth}{../code/Synthetic/figures/evaluateCrossValidationResults_Synthetic_LaplaceNoise.py_SimulateSynthetic_Parameterized_OtherNoiseLevels_Grid_VarySize_LaplaceNoise.py_180_8_2345_N40000_STEEPPERIODIC_UNIFORM.txt_RelativeLF.pdf}{3.0cm}{NLL p=8 40k} \\
\safegraphic{0.14\textwidth}{../code/Synthetic/figures/CounterfactualModel_VIZ_ByNoiseMagnGauge_Components.py_2345_UNIFORM_STEEPPERIODIC_2_0_10.0_180.pdf}{2.2cm}{gt uniform periodic} &
\safegraphic{0.18\textwidth}{../code/Synthetic/figures/RunSynthetic_FreePrior_CosineLoss_OnSim_LaplaceNoise_VIZ_OnlyModel_OtherNoiseLevels.py_SimulateSynthetic_Parameterized_OtherNoiseLevels_Grid_VarySize_LaplaceNoise.py_180_2_2345_N40000_UNIFORM_STEEPPERIODIC.txt_2_0_10.0_180.pdf}{3.0cm}{laplace fit p=2 40k} &
\safegraphic{0.10\textwidth}{../code/Synthetic/figures/evaluateCrossValidationResults_Synthetic_LaplaceNoise.py_SimulateSynthetic_Parameterized_OtherNoiseLevels_Grid_VarySize_LaplaceNoise.py_180_2_2345_N40000_UNIFORM_STEEPPERIODIC.txt_RelativeLF.pdf}{3.0cm}{NLL p=2 40k} &
\safegraphic{0.18\textwidth}{../code/Synthetic/figures/RunSynthetic_FreePrior_CosineLoss_OnSim_LaplaceNoise_VIZ_OnlyModel_OtherNoiseLevels.py_SimulateSynthetic_Parameterized_OtherNoiseLevels_Grid_VarySize_LaplaceNoise.py_180_8_2345_N40000_UNIFORM_STEEPPERIODIC.txt_8_0_10.0_180.pdf}{3.0cm}{laplace fit p=8 40k} &
\safegraphic{0.10\textwidth}{../code/Synthetic/figures/evaluateCrossValidationResults_Synthetic_LaplaceNoise.py_SimulateSynthetic_Parameterized_OtherNoiseLevels_Grid_VarySize_LaplaceNoise.py_180_8_2345_N40000_UNIFORM_STEEPPERIODIC.txt_RelativeLF.pdf}{3.0cm}{NLL p=8 40k} \\
\safegraphic{0.14\textwidth}{../code/Synthetic/figures/CounterfactualModel_VIZ_ByNoiseMagnGauge_Components.py_2345_STEEPSHIFTED_STEEPPERIODIC_2_0_10.0_180.pdf}{2.2cm}{gt shifted periodic} &
\safegraphic{0.18\textwidth}{../code/Synthetic/figures/RunSynthetic_FreePrior_CosineLoss_OnSim_LaplaceNoise_VIZ_OnlyModel_OtherNoiseLevels.py_SimulateSynthetic_Parameterized_OtherNoiseLevels_Grid_VarySize_LaplaceNoise.py_180_2_2345_N40000_STEEPSHIFTED_STEEPPERIODIC.txt_2_0_10.0_180.pdf}{3.0cm}{laplace fit p=2 40k} &
\safegraphic{0.10\textwidth}{../code/Synthetic/figures/evaluateCrossValidationResults_Synthetic_LaplaceNoise.py_SimulateSynthetic_Parameterized_OtherNoiseLevels_Grid_VarySize_LaplaceNoise.py_180_2_2345_N40000_STEEPSHIFTED_STEEPPERIODIC.txt_RelativeLF.pdf}{3.0cm}{NLL p=2 40k} &
\safegraphic{0.18\textwidth}{../code/Synthetic/figures/RunSynthetic_FreePrior_CosineLoss_OnSim_LaplaceNoise_VIZ_OnlyModel_OtherNoiseLevels.py_SimulateSynthetic_Parameterized_OtherNoiseLevels_Grid_VarySize_LaplaceNoise.py_180_8_2345_N40000_STEEPSHIFTED_STEEPPERIODIC.txt_8_0_10.0_180.pdf}{3.0cm}{laplace fit p=8 40k} &
\safegraphic{0.10\textwidth}{../code/Synthetic/figures/evaluateCrossValidationResults_Synthetic_LaplaceNoise.py_SimulateSynthetic_Parameterized_OtherNoiseLevels_Grid_VarySize_LaplaceNoise.py_180_8_2345_N40000_STEEPSHIFTED_STEEPPERIODIC.txt_RelativeLF.pdf}{3.0cm}{NLL p=8 40k}
\end{tabular}
\caption{Applying our fitting procedure (but assuming Laplace-like sensory noise) to data generated under Laplace-like sensory noise.
Note that, for the model fits, the resources $\sqrt{\mathcal{J}(\theta)}$ are computed as $F'(\theta)/\sigma$ in keeping with our other (von Mises/Gaussian) fits; results may slightly differ if computed using the Fisher Information for the Laplace distribution.
}
\label{fig:si-alt-laplace-on-laplace}
\end{figure}


\begin{comment}
python3 RunSynthetic_FreePrior_CosineLoss_OnSim_LaplaceNoise_VIZ_OnlyModel_OtherNoiseLevels.py 2 0 10.0 180 "SimulateSynthetic_Parameterized_OtherNoiseLevels_Grid_VarySize_LaplaceNoise.py_180_2_2345_N10000_STEEPPERIODIC_STEEPPERIODIC.txt"
python3 RunSynthetic_FreePrior_CosineLoss_OnSim_LaplaceNoise_VIZ_OnlyModel_OtherNoiseLevels.py 2 0 10.0 180 "SimulateSynthetic_Parameterized_OtherNoiseLevels_Grid_VarySize_LaplaceNoise.py_180_2_2345_N10000_STEEPPERIODIC_UNIFORM.txt"
python3 RunSynthetic_FreePrior_CosineLoss_OnSim_LaplaceNoise_VIZ_OnlyModel_OtherNoiseLevels.py 2 0 10.0 180 "SimulateSynthetic_Parameterized_OtherNoiseLevels_Grid_VarySize_LaplaceNoise.py_180_2_2345_N10000_UNIFORM_STEEPPERIODIC.txt"
python3 RunSynthetic_FreePrior_CosineLoss_OnSim_LaplaceNoise_VIZ_OnlyModel_OtherNoiseLevels.py 2 0 10.0 180 "SimulateSynthetic_Parameterized_OtherNoiseLevels_Grid_VarySize_LaplaceNoise.py_180_2_2345_N10000_STEEPSHIFTED_STEEPPERIODIC.txt"
python3 RunSynthetic_FreePrior_CosineLoss_OnSim_LaplaceNoise_VIZ_OnlyModel_OtherNoiseLevels.py 8 0 10.0 180 "SimulateSynthetic_Parameterized_OtherNoiseLevels_Grid_VarySize_LaplaceNoise.py_180_8_2345_N10000_STEEPPERIODIC_STEEPPERIODIC.txt"
python3 RunSynthetic_FreePrior_CosineLoss_OnSim_LaplaceNoise_VIZ_OnlyModel_OtherNoiseLevels.py 8 0 10.0 180 "SimulateSynthetic_Parameterized_OtherNoiseLevels_Grid_VarySize_LaplaceNoise.py_180_8_2345_N10000_STEEPPERIODIC_UNIFORM.txt"
python3 RunSynthetic_FreePrior_CosineLoss_OnSim_LaplaceNoise_VIZ_OnlyModel_OtherNoiseLevels.py 8 0 10.0 180 "SimulateSynthetic_Parameterized_OtherNoiseLevels_Grid_VarySize_LaplaceNoise.py_180_8_2345_N10000_UNIFORM_STEEPPERIODIC.txt"
python3 RunSynthetic_FreePrior_CosineLoss_OnSim_LaplaceNoise_VIZ_OnlyModel_OtherNoiseLevels.py 8 0 10.0 180 "SimulateSynthetic_Parameterized_OtherNoiseLevels_Grid_VarySize_LaplaceNoise.py_180_8_2345_N10000_STEEPSHIFTED_STEEPPERIODIC.txt"

python3 evaluateCrossValidationResults_Synthetic_LaplaceNoise.py "SimulateSynthetic_Parameterized_OtherNoiseLevels_Grid_VarySize_LaplaceNoise.py_180_2_2345_N10000_STEEPPERIODIC_STEEPPERIODIC.txt"
python3 evaluateCrossValidationResults_Synthetic_LaplaceNoise.py "SimulateSynthetic_Parameterized_OtherNoiseLevels_Grid_VarySize_LaplaceNoise.py_180_2_2345_N10000_STEEPPERIODIC_UNIFORM.txt"
python3 evaluateCrossValidationResults_Synthetic_LaplaceNoise.py "SimulateSynthetic_Parameterized_OtherNoiseLevels_Grid_VarySize_LaplaceNoise.py_180_2_2345_N10000_UNIFORM_STEEPPERIODIC.txt"
python3 evaluateCrossValidationResults_Synthetic_LaplaceNoise.py "SimulateSynthetic_Parameterized_OtherNoiseLevels_Grid_VarySize_LaplaceNoise.py_180_2_2345_N10000_STEEPSHIFTED_STEEPPERIODIC.txt"
python3 evaluateCrossValidationResults_Synthetic_LaplaceNoise.py "SimulateSynthetic_Parameterized_OtherNoiseLevels_Grid_VarySize_LaplaceNoise.py_180_8_2345_N10000_STEEPPERIODIC_STEEPPERIODIC.txt"
python3 evaluateCrossValidationResults_Synthetic_LaplaceNoise.py "SimulateSynthetic_Parameterized_OtherNoiseLevels_Grid_VarySize_LaplaceNoise.py_180_8_2345_N10000_STEEPPERIODIC_UNIFORM.txt"
python3 evaluateCrossValidationResults_Synthetic_LaplaceNoise.py "SimulateSynthetic_Parameterized_OtherNoiseLevels_Grid_VarySize_LaplaceNoise.py_180_8_2345_N10000_UNIFORM_STEEPPERIODIC.txt"
python3 evaluateCrossValidationResults_Synthetic_LaplaceNoise.py "SimulateSynthetic_Parameterized_OtherNoiseLevels_Grid_VarySize_LaplaceNoise.py_180_8_2345_N10000_STEEPSHIFTED_STEEPPERIODIC.txt"

python3 RunSynthetic_FreePrior_CosineLoss_OnSim_LaplaceNoise_VIZ_OnlyModel_OtherNoiseLevels.py 2 0 10.0 180 "SimulateSynthetic_Parameterized_OtherNoiseLevels_Grid_VarySize_LaplaceNoise.py_180_2_2345_N40000_STEEPPERIODIC_STEEPPERIODIC.txt"
python3 RunSynthetic_FreePrior_CosineLoss_OnSim_LaplaceNoise_VIZ_OnlyModel_OtherNoiseLevels.py 2 0 10.0 180 "SimulateSynthetic_Parameterized_OtherNoiseLevels_Grid_VarySize_LaplaceNoise.py_180_2_2345_N40000_STEEPPERIODIC_UNIFORM.txt"
python3 RunSynthetic_FreePrior_CosineLoss_OnSim_LaplaceNoise_VIZ_OnlyModel_OtherNoiseLevels.py 2 0 10.0 180 "SimulateSynthetic_Parameterized_OtherNoiseLevels_Grid_VarySize_LaplaceNoise.py_180_2_2345_N40000_UNIFORM_STEEPPERIODIC.txt"
python3 RunSynthetic_FreePrior_CosineLoss_OnSim_LaplaceNoise_VIZ_OnlyModel_OtherNoiseLevels.py 2 0 10.0 180 "SimulateSynthetic_Parameterized_OtherNoiseLevels_Grid_VarySize_LaplaceNoise.py_180_2_2345_N40000_STEEPSHIFTED_STEEPPERIODIC.txt"
python3 RunSynthetic_FreePrior_CosineLoss_OnSim_LaplaceNoise_VIZ_OnlyModel_OtherNoiseLevels.py 8 0 10.0 180 "SimulateSynthetic_Parameterized_OtherNoiseLevels_Grid_VarySize_LaplaceNoise.py_180_8_2345_N40000_STEEPPERIODIC_STEEPPERIODIC.txt"
python3 RunSynthetic_FreePrior_CosineLoss_OnSim_LaplaceNoise_VIZ_OnlyModel_OtherNoiseLevels.py 8 0 10.0 180 "SimulateSynthetic_Parameterized_OtherNoiseLevels_Grid_VarySize_LaplaceNoise.py_180_8_2345_N40000_STEEPPERIODIC_UNIFORM.txt"
python3 RunSynthetic_FreePrior_CosineLoss_OnSim_LaplaceNoise_VIZ_OnlyModel_OtherNoiseLevels.py 8 0 10.0 180 "SimulateSynthetic_Parameterized_OtherNoiseLevels_Grid_VarySize_LaplaceNoise.py_180_8_2345_N40000_UNIFORM_STEEPPERIODIC.txt"
python3 RunSynthetic_FreePrior_CosineLoss_OnSim_LaplaceNoise_VIZ_OnlyModel_OtherNoiseLevels.py 8 0 10.0 180 "SimulateSynthetic_Parameterized_OtherNoiseLevels_Grid_VarySize_LaplaceNoise.py_180_8_2345_N40000_STEEPSHIFTED_STEEPPERIODIC.txt"

python3 evaluateCrossValidationResults_Synthetic_LaplaceNoise.py "SimulateSynthetic_Parameterized_OtherNoiseLevels_Grid_VarySize_LaplaceNoise.py_180_2_2345_N40000_STEEPPERIODIC_STEEPPERIODIC.txt"
python3 evaluateCrossValidationResults_Synthetic_LaplaceNoise.py "SimulateSynthetic_Parameterized_OtherNoiseLevels_Grid_VarySize_LaplaceNoise.py_180_2_2345_N40000_STEEPPERIODIC_UNIFORM.txt"
python3 evaluateCrossValidationResults_Synthetic_LaplaceNoise.py "SimulateSynthetic_Parameterized_OtherNoiseLevels_Grid_VarySize_LaplaceNoise.py_180_2_2345_N40000_UNIFORM_STEEPPERIODIC.txt"
python3 evaluateCrossValidationResults_Synthetic_LaplaceNoise.py "SimulateSynthetic_Parameterized_OtherNoiseLevels_Grid_VarySize_LaplaceNoise.py_180_2_2345_N40000_STEEPSHIFTED_STEEPPERIODIC.txt"
python3 evaluateCrossValidationResults_Synthetic_LaplaceNoise.py "SimulateSynthetic_Parameterized_OtherNoiseLevels_Grid_VarySize_LaplaceNoise.py_180_8_2345_N40000_STEEPPERIODIC_STEEPPERIODIC.txt"
python3 evaluateCrossValidationResults_Synthetic_LaplaceNoise.py "SimulateSynthetic_Parameterized_OtherNoiseLevels_Grid_VarySize_LaplaceNoise.py_180_8_2345_N40000_STEEPPERIODIC_UNIFORM.txt"
python3 evaluateCrossValidationResults_Synthetic_LaplaceNoise.py "SimulateSynthetic_Parameterized_OtherNoiseLevels_Grid_VarySize_LaplaceNoise.py_180_8_2345_N40000_UNIFORM_STEEPPERIODIC.txt"
python3 evaluateCrossValidationResults_Synthetic_LaplaceNoise.py "SimulateSynthetic_Parameterized_OtherNoiseLevels_Grid_VarySize_LaplaceNoise.py_180_8_2345_N40000_STEEPSHIFTED_STEEPPERIODIC.txt"
\end{comment}


\clearpage
\newpage


\subsection{Beyond $L_p$ Losses}

Our theoretical results and simulations cover the $L_p$ family of loss functions.
Here, we discuss the possibility of expanding our theory to other loss functions.

\paragraph*{Symmetric Losses}
A natural constraint on loss functions is that they should be symmetric, i.e., $\ell(x,y) = \ell(y,x)$, and that they should be isotropic, i.e., $\ell(x,y)$ only depends on $x-y$.
Here, we consider a general isotropic and symmetric smooth loss $\ell(x,y)$. We assume that $\ell(x,y)$ takes a unique minimum at $x=y$; we may assume $\ell(x,y) = 0$ in this case.
Then, using a Taylor expansion around $x-y=0$:
\begin{equation}
\ell(x,y) = \alpha_1 (x-y)^p + \alpha_2 (x-y)^{p+2} + \mathcal{O}(|x-y|^{p+3})
\end{equation}
for some even $p>0$ and some $\alpha_1 > 0$.
When noise is small, this is dominated by the lowest-order term, i.e., $\ell(x,y)$ behaves like $L_p$ loss, up to the scaling factor $\alpha_1$, which does not impact the minimizer:
\begin{equation}\label{eq:smooth-loss-lowest-order}
\ell(x,y) = \alpha_1 (x-y)^p + \mathcal{O}(|x-y|^{p+2})
\end{equation}
One may thus conjecture that, in this regime, recovering the order $p$ is possible under similar circumstances as recovery of the exponent of $L_p$ loss.
If one desires to recover the full shape of $\ell$, further arguments are likely needed. 
One plausible approach here may be to attempt to recover the higher-order Taylor expansion coefficients of $\ell$ (around $x-y=0$) by observing the bias under a sufficiently large number of distinct sensory noise levels; we leave exploration of this idea to future work.


As an illustration of these considerations, we consider the Inverted Gaussian loss function \citep{Krding2004TheLF, Acerbi2014AFF}:
\begin{equation}\label{eq:acerbi-loss}
\ell(x,y) = 1 - \exp(-\frac{(x-y)^2}{2\sigma^2})
\end{equation}
This loss function was proposed by \citet{Krding2004TheLF} as a parametric form for a loss function that increases less than $L_2$ loss for large errors.
In their identifiability simulations, \citet{Acerbi2014AFF} used such functions with various choices of $\sigma$.\footnote{In their theoretical calculations, \citet{Acerbi2014AFF} considered a more general family consisting of linear combinations of pairs of such functions. In their simulation studies (their Sections 4.1 and 4.2), they considered the Inverted Gaussian loss with different $\sigma$ parameters.}
We show the loss function at three values of $\sigma^2$ in Figure~\ref{fig:acerbi-loss-family-preview-remington}; it behaves similarly to scaled versions of $L_2$ loss when errors are close to 0, though the behavior at larger errors is different.

Eq.~\ref{eq:acerbi-loss} defines a symmetric and smooth loss function, with a unique minimum at $x-y=0$.
A Taylor expansion around $x=y$ shows (compare the general form in (\ref{eq:smooth-loss-lowest-order})):
\begin{equation}
\ell(x,y) = \frac{(x-y)^2}{2\sigma^2} + \mathcal{O}((x-y)^4)
\end{equation}
Thus, the loss function, in the small-noise regime, behaves similarly to $L_2$ loss.

We simulated data from this loss function in a way analogous to our interval space simulations~\ref{sec:further-interval}, but with this loss function, at varying values of $\sigma^2$.
We then applied our fitting procedure (assuming $L_p$ losses) to these datasets. We show results in Figure~\ref{fig:si-alt-acerbi-remington-preview}.
Result show that fit is best at $p=2$, and that encoding and prior are often accurately recovered when fitting using $L_2$ loss, despite the mismatch between the true loss and the fitted $L_p$ family.




\begin{figure}[h]
\centering
\includegraphics[width=0.5\textwidth]{figures/acerbi_loss_shapes_interval.pdf}
\caption{Inverted Gaussian loss functions based on \citet{Krding2004TheLF,Acerbi2014AFF}, for different values of $\sigma$ (Equation~\ref{eq:acerbi-loss}).}
\label{fig:acerbi-loss-family-preview-remington}
\end{figure}

\begin{figure}[p]
\centering
\scriptsize
\setlength{\tabcolsep}{2pt}
\newcommand{\groundtruthacerbiremington}[1]{%
  \adjustbox{valign=m,max totalsize={0.1575\textwidth}{2.2cm},keepaspectratio}{\includegraphics{#1}}%
}
\begin{tabular}{>{\centering\arraybackslash}m{0.1575\textwidth}@{\hspace{0.010\textwidth}}>{\centering\arraybackslash}m{0.1575\textwidth}>{\centering\arraybackslash}m{0.07875\textwidth}@{\hspace{0.045\textwidth}}>{\centering\arraybackslash}m{0.1575\textwidth}>{\centering\arraybackslash}m{0.07875\textwidth}@{\hspace{0.045\textwidth}}>{\centering\arraybackslash}m{0.1575\textwidth}>{\centering\arraybackslash}m{0.07875\textwidth}}
\textbf{Model} &
\multicolumn{2}{c}{\textbf{{$\sigma = 0.1$}}} &
\multicolumn{2}{c}{\textbf{{$\sigma = 0.5$}}} &
\multicolumn{2}{c}{\textbf{{$\sigma = 1.0$}}} \\
& \textbf{Fit} & \textbf{NLL} & \textbf{Fit} & \textbf{NLL} & \textbf{Fit} & \textbf{NLL} \\
\groundtruthacerbiremington{figures/CounterfactualModel_Remington_VIZ_Components.py_4567_UNIFORM_UNIMODAL2_2_0_10.0_400.pdf} &
\safegraphic{0.1575\textwidth}{../code/Remington/figures/RunSynthetic_DenseRemington_FreeEncoding_OnSim_OtherNoiseLevels_VarySize_AcerbiLoss_VIZ_OnlyModel.py_SimAcer.py_400_2_4567_N10000_UNIFORM_UNIMODAL2_pi_1_mean1_0.0_mean2_0.0_sigma21_0.1_sigma22_0.1.txt_2_0_10.0_400.pdf}{2.2cm}{fit} &
\safegraphic{0.07875\textwidth}{../code/Remington/figures/evaluateCrossValidationResults_Synthetic_DenseRemington_AcerbiLoss.py_SimAcer.py_400_2_4567_N10000_UNIFORM_UNIMODAL2_pi_1_mean1_0.0_mean2_0.0_sigma21_0.1_sigma22_0.1.txt_RelativeLF.pdf}{2.2cm}{NLL} &
\safegraphic{0.1575\textwidth}{../code/Remington/figures/RunSynthetic_DenseRemington_FreeEncoding_OnSim_OtherNoiseLevels_VarySize_AcerbiLoss_VIZ_OnlyModel.py_SimAcer.py_400_2_4567_N10000_UNIFORM_UNIMODAL2_pi_1_mean1_0.0_mean2_0.0_sigma21_0.5_sigma22_0.1.txt_2_0_10.0_400.pdf}{2.2cm}{fit} &
\safegraphic{0.07875\textwidth}{../code/Remington/figures/evaluateCrossValidationResults_Synthetic_DenseRemington_AcerbiLoss.py_SimAcer.py_400_2_4567_N10000_UNIFORM_UNIMODAL2_pi_1_mean1_0.0_mean2_0.0_sigma21_0.5_sigma22_0.1.txt_RelativeLF.pdf}{2.2cm}{NLL} &
\safegraphic{0.1575\textwidth}{../code/Remington/figures/RunSynthetic_DenseRemington_FreeEncoding_OnSim_OtherNoiseLevels_VarySize_AcerbiLoss_VIZ_OnlyModel.py_SimAcer.py_400_2_4567_N10000_UNIFORM_UNIMODAL2_pi_1_mean1_0.0_mean2_0.0_sigma21_1.0_sigma22_0.1.txt_2_0_10.0_400.pdf}{2.2cm}{fit} &
\safegraphic{0.07875\textwidth}{../code/Remington/figures/evaluateCrossValidationResults_Synthetic_DenseRemington_AcerbiLoss.py_SimAcer.py_400_2_4567_N10000_UNIFORM_UNIMODAL2_pi_1_mean1_0.0_mean2_0.0_sigma21_1.0_sigma22_0.1.txt_RelativeLF.pdf}{2.2cm}{NLL} \\
\groundtruthacerbiremington{figures/CounterfactualModel_Remington_VIZ_Components.py_4567_UNIMODAL2_UNIMODAL2_2_0_10.0_400.pdf} &
\safegraphic{0.1575\textwidth}{../code/Remington/figures/RunSynthetic_DenseRemington_FreeEncoding_OnSim_OtherNoiseLevels_VarySize_AcerbiLoss_VIZ_OnlyModel.py_SimAcer.py_400_2_4567_N10000_UNIMODAL2_UNIMODAL2_pi_1_mean1_0.0_mean2_0.0_sigma21_0.1_sigma22_0.1.txt_2_0_10.0_400.pdf}{2.2cm}{fit} &
\safegraphic{0.07875\textwidth}{../code/Remington/figures/evaluateCrossValidationResults_Synthetic_DenseRemington_AcerbiLoss.py_SimAcer.py_400_2_4567_N10000_UNIMODAL2_UNIMODAL2_pi_1_mean1_0.0_mean2_0.0_sigma21_0.1_sigma22_0.1.txt_RelativeLF.pdf}{2.2cm}{NLL} &
\safegraphic{0.1575\textwidth}{../code/Remington/figures/RunSynthetic_DenseRemington_FreeEncoding_OnSim_OtherNoiseLevels_VarySize_AcerbiLoss_VIZ_OnlyModel.py_SimAcer.py_400_2_4567_N10000_UNIMODAL2_UNIMODAL2_pi_1_mean1_0.0_mean2_0.0_sigma21_0.5_sigma22_0.1.txt_2_0_10.0_400.pdf}{2.2cm}{fit} &
\safegraphic{0.07875\textwidth}{../code/Remington/figures/evaluateCrossValidationResults_Synthetic_DenseRemington_AcerbiLoss.py_SimAcer.py_400_2_4567_N10000_UNIMODAL2_UNIMODAL2_pi_1_mean1_0.0_mean2_0.0_sigma21_0.5_sigma22_0.1.txt_RelativeLF.pdf}{2.2cm}{NLL} &
\safegraphic{0.1575\textwidth}{../code/Remington/figures/RunSynthetic_DenseRemington_FreeEncoding_OnSim_OtherNoiseLevels_VarySize_AcerbiLoss_VIZ_OnlyModel.py_SimAcer.py_400_2_4567_N10000_UNIMODAL2_UNIMODAL2_pi_1_mean1_0.0_mean2_0.0_sigma21_1.0_sigma22_0.1.txt_2_0_10.0_400.pdf}{2.2cm}{fit} &
\safegraphic{0.07875\textwidth}{../code/Remington/figures/evaluateCrossValidationResults_Synthetic_DenseRemington_AcerbiLoss.py_SimAcer.py_400_2_4567_N10000_UNIMODAL2_UNIMODAL2_pi_1_mean1_0.0_mean2_0.0_sigma21_1.0_sigma22_0.1.txt_RelativeLF.pdf}{2.2cm}{NLL} \\
\groundtruthacerbiremington{figures/CounterfactualModel_Remington_VIZ_Components.py_4567_UNIFORM_BIMODAL2_2_0_10.0_400.pdf} &
\safegraphic{0.1575\textwidth}{../code/Remington/figures/RunSynthetic_DenseRemington_FreeEncoding_OnSim_OtherNoiseLevels_VarySize_AcerbiLoss_VIZ_OnlyModel.py_SimAcer.py_400_2_4567_N10000_UNIFORM_BIMODAL2_pi_1_mean1_0.0_mean2_0.0_sigma21_0.1_sigma22_0.1.txt_2_0_10.0_400.pdf}{2.2cm}{fit} &
\safegraphic{0.07875\textwidth}{../code/Remington/figures/evaluateCrossValidationResults_Synthetic_DenseRemington_AcerbiLoss.py_SimAcer.py_400_2_4567_N10000_UNIFORM_BIMODAL2_pi_1_mean1_0.0_mean2_0.0_sigma21_0.1_sigma22_0.1.txt_RelativeLF.pdf}{2.2cm}{NLL} &
\safegraphic{0.1575\textwidth}{../code/Remington/figures/RunSynthetic_DenseRemington_FreeEncoding_OnSim_OtherNoiseLevels_VarySize_AcerbiLoss_VIZ_OnlyModel.py_SimAcer.py_400_2_4567_N10000_UNIFORM_BIMODAL2_pi_1_mean1_0.0_mean2_0.0_sigma21_0.5_sigma22_0.1.txt_2_0_10.0_400.pdf}{2.2cm}{fit} &
\safegraphic{0.07875\textwidth}{../code/Remington/figures/evaluateCrossValidationResults_Synthetic_DenseRemington_AcerbiLoss.py_SimAcer.py_400_2_4567_N10000_UNIFORM_BIMODAL2_pi_1_mean1_0.0_mean2_0.0_sigma21_0.5_sigma22_0.1.txt_RelativeLF.pdf}{2.2cm}{NLL} &
\safegraphic{0.1575\textwidth}{../code/Remington/figures/RunSynthetic_DenseRemington_FreeEncoding_OnSim_OtherNoiseLevels_VarySize_AcerbiLoss_VIZ_OnlyModel.py_SimAcer.py_400_2_4567_N10000_UNIFORM_BIMODAL2_pi_1_mean1_0.0_mean2_0.0_sigma21_1.0_sigma22_0.1.txt_2_0_10.0_400.pdf}{2.2cm}{fit} &
\safegraphic{0.07875\textwidth}{../code/Remington/figures/evaluateCrossValidationResults_Synthetic_DenseRemington_AcerbiLoss.py_SimAcer.py_400_2_4567_N10000_UNIFORM_BIMODAL2_pi_1_mean1_0.0_mean2_0.0_sigma21_1.0_sigma22_0.1.txt_RelativeLF.pdf}{2.2cm}{NLL} \\
\groundtruthacerbiremington{figures/CounterfactualModel_Remington_VIZ_Components.py_4567_UNIMODAL2_UNIFORM_2_0_10.0_400.pdf} &
\safegraphic{0.1575\textwidth}{../code/Remington/figures/RunSynthetic_DenseRemington_FreeEncoding_OnSim_OtherNoiseLevels_VarySize_AcerbiLoss_VIZ_OnlyModel.py_SimAcer.py_400_2_4567_N10000_UNIMODAL2_UNIFORM_pi_1_mean1_0.0_mean2_0.0_sigma21_0.1_sigma22_0.1.txt_2_0_10.0_400.pdf}{2.2cm}{fit} &
\safegraphic{0.07875\textwidth}{../code/Remington/figures/evaluateCrossValidationResults_Synthetic_DenseRemington_AcerbiLoss.py_SimAcer.py_400_2_4567_N10000_UNIMODAL2_UNIFORM_pi_1_mean1_0.0_mean2_0.0_sigma21_0.1_sigma22_0.1.txt_RelativeLF.pdf}{2.2cm}{NLL} &
\safegraphic{0.1575\textwidth}{../code/Remington/figures/RunSynthetic_DenseRemington_FreeEncoding_OnSim_OtherNoiseLevels_VarySize_AcerbiLoss_VIZ_OnlyModel.py_SimAcer.py_400_2_4567_N10000_UNIMODAL2_UNIFORM_pi_1_mean1_0.0_mean2_0.0_sigma21_0.5_sigma22_0.1.txt_2_0_10.0_400.pdf}{2.2cm}{fit} &
\safegraphic{0.07875\textwidth}{../code/Remington/figures/evaluateCrossValidationResults_Synthetic_DenseRemington_AcerbiLoss.py_SimAcer.py_400_2_4567_N10000_UNIMODAL2_UNIFORM_pi_1_mean1_0.0_mean2_0.0_sigma21_0.5_sigma22_0.1.txt_RelativeLF.pdf}{2.2cm}{NLL} &
\safegraphic{0.1575\textwidth}{../code/Remington/figures/RunSynthetic_DenseRemington_FreeEncoding_OnSim_OtherNoiseLevels_VarySize_AcerbiLoss_VIZ_OnlyModel.py_SimAcer.py_400_2_4567_N10000_UNIMODAL2_UNIFORM_pi_1_mean1_0.0_mean2_0.0_sigma21_1.0_sigma22_0.1.txt_2_0_10.0_400.pdf}{2.2cm}{fit} &
\safegraphic{0.07875\textwidth}{../code/Remington/figures/evaluateCrossValidationResults_Synthetic_DenseRemington_AcerbiLoss.py_SimAcer.py_400_2_4567_N10000_UNIMODAL2_UNIFORM_pi_1_mean1_0.0_mean2_0.0_sigma21_1.0_sigma22_0.1.txt_RelativeLF.pdf}{2.2cm}{NLL} \\
\end{tabular}
\caption{Here, we simulated data from four different prior-encoding pairs (as in Section~\ref{sec:further-interval}; four noise levels at 10K trials each) using three different versions of the inverted Gaussian loss function (Figure~\ref{fig:acerbi-loss-family-preview-remington} and Equation~\ref{eq:acerbi-loss}; $\sigma$ controls the shape of the loss function), and then applied our standard fitting procedure.
Throughout, best fit is achieved at $p=2$ or close to it, in line with theoretical expectations.
For each dataset, we show the fit at $L_2$ loss.
The shape of the encoding is adequately recovered despite the mismatch between fitted ($L_2$) and ground truth (inverted Gaussian loss), across levels of $\sigma$.
The prior is accurately recovered at higher $\sigma$; at low $\sigma$ the shape is accurate except in the third row.
This suggests that the encoding can be accurately recovered even if the loss function is misspecified; the prior may have some sensitivity to loss misspecification, but is still often accurately recovered.
}
\label{fig:si-alt-acerbi-remington-preview}
\end{figure}

\begin{comment}
for sigma in 0.1 0.2 0.4 0.5 1.0; do
  for fit in \
    "SimulateSynthetic2_DenseRemington_OtherNoiseLevels_VarySize_AcerbiLoss.py_400_2_4567_N10000_UNIFORM_UNIMODAL2_pi_1_mean1_0.0_mean2_0.0_sigma21_${sigma}_sigma22_0.1.txt" \
    "SimulateSynthetic2_DenseRemington_OtherNoiseLevels_VarySize_AcerbiLoss.py_400_2_4567_N10000_UNIMODAL2_UNIMODAL2_pi_1_mean1_0.0_mean2_0.0_sigma21_${sigma}_sigma22_0.1.txt" \
    "SimulateSynthetic2_DenseRemington_OtherNoiseLevels_VarySize_AcerbiLoss.py_400_2_4567_N10000_UNIFORM_BIMODAL2_pi_1_mean1_0.0_mean2_0.0_sigma21_${sigma}_sigma22_0.1.txt" \
    "SimulateSynthetic2_DenseRemington_OtherNoiseLevels_VarySize_AcerbiLoss.py_400_2_4567_N10000_UNIMODAL2_UNIFORM_pi_1_mean1_0.0_mean2_0.0_sigma21_${sigma}_sigma22_0.1.txt"
  do
    python3 evaluateCrossValidationResults_Synthetic_DenseRemington_AcerbiLoss.py "$fit"
    python3 RunSynthetic_DenseRemington_FreeEncoding_Zero_OnSim_OtherNoiseLevels_VarySize_AcerbiLoss_Round2_VIZ_OnlyModel.py 0 0 10.0 400 "$fit"
    python3 RunSynthetic_DenseRemington_FreeEncoding_L1_OnSim_OtherNoiseLevels_VarySize_AcerbiLoss_Round2_VIZ_OnlyModel.py 1 0 10.0 400 "$fit"
    python3 RunSynthetic_DenseRemington_FreeEncoding_OnSim_OtherNoiseLevels_VarySize_AcerbiLoss_VIZ_OnlyModel.py 2 0 10.0 400 "$fit"
  done
done
\end{comment}


\paragraph*{Asymmetric Losses}
Even more general loss functions might abandon the assumption symmetry; for instance, by differentially penalizing over- vs under-estimation \citep[e.g.][]{Mamassian2008OverconfidenceIA,Hudson2008OptimalCF}.
Such loss functions may lead to qualitatively different patterns of bias (such as a systematic bias towards over- or under-estimation); thus, the asymmetry would need to be distinguished from a skew in the prior. 
A potential approach to expanding identifiability theory to such loss functions may be to again consider a Taylor expansion, which in the case of asymmetric losses would include odd powers of $x-y$ (for even $p$)\footnote{$p$ must be even to ensure that the loss takes a minimum at $x=y$, but higher odd orders can have nonzero coefficients.}, e.g.:
\begin{equation}
\ell(x,y) = \alpha_1 \cdot (x-y)^p + \alpha_2 \cdot (x-y)^{p+1} + \mathcal{O}(|x-y|^{p+2})
\end{equation}
and investigate how the coefficients $\alpha_1, \alpha_2, \dots$ might be recovered from the expansion of the bias and variability (Lemma~\ref{lemma:expansion}), possibly requiring a higher order ($t^3$, $t^4$ etc), and possibly under a sufficiently large number of distinct noise levels.
We leave this investigation to future work.

\clearpage
\newpage



\section{Details for Fitting Procedure}
\label{sec:tutorial}

Here, we provide a self-contained tutorial for the model and fitting procedure, with the goal of making it easier for readers to understand the model and reproduce the results.
The fitting procedure is from \citet{hahn2024unifying}, and we refer to the formal description there.
Here, we intend to give a user-friendly exposition to readers interested in applying the fitting procedure to psychophysical data.

The codebase provides instructions for running on a simulated dataset\footnote{\url{https://github.com/m-hahn/identifiability-bayesian-models/tree/main/code}}, on the original dataset by \citet{deGardelle2010AnOI}\footnote{\url{https://github.com/m-hahn/identifiability-bayesian-models/tree/main/code/Gardelle}} and on the dataset by \citet{Remington2018LateBI}\footnote{\url{https://gitlab.com/m-hahn/unifying-theory-biases/-/tree/main/code/Remington} in the codebase from \citet{hahn2024unifying}.}, including for visualizing the resulting fits.

We first discuss which aspects are specified by the user, and which parts are recovered by the fitting procedure.
The user specifies the following parts:
\begin{enumerate}
\item The stimulus space $\mathcal{X}$; it may be a finite interval or a circle.
The  sensory space $\mathcal{Y}$ has the same topology; in practice, it can be chosen to be $[0,1]$ (interval) or $[0,2\pi)$ (circular).
\item The exponent $p \geq 0$ of the $L_p$ loss function. The fitting procedure admits any nonnegative integer; we use $p=0,1,2,4,6,8,\dots$.
\item The number $L$ of sensory noise levels.
\end{enumerate}
The behavioral data is expected to consist of trials that each consist of
\begin{enumerate}
\item the true stimulus $\theta \in \mathcal{X}$,
\item the observed response $x \in \mathcal{X}$,
\item an index for the sensory noise level, $l \in \{1, \dots, L\}$ (e.g., an index for a condition under a manipulation of exposure duration, contrast, or another manipulation affecting the magnitude of sensory noise).
\end{enumerate}
Given these, the fitting procedure determines
\begin{enumerate}
	\item A prior $\Prior(\theta)$, which is a probability density on $\mathcal{X}$.
	\item An encoding function $F(\theta) : \mathcal{X} \rightarrow \mathcal{Y}$; equivalently, a function $F'(\theta) : \mathcal{X} \mapsto (0,\infty)$ describing the slope of $F(\theta)$.
	\item Sensory noise magnitudes  $\sigma_1, \dots, \sigma_L > 0$, which specify the magnitude of sensory noise in the $L$ different noise level conditions
        \item The variance  $\rho^2$ of motor noise. 
	\item $\xi \in [0,1)$, the guessing rate.
\end{enumerate}
There are several possible extensions, such as to stimulus (external) noise (Section~\ref{sec:stimulus-noise}), to subject-specific adjustments (SI Appendix S2.1.4 in \citet{hahn2024unifying}), and to encoding functions varying with noise conditions (Section~\ref{sec:vary-with-noise}).

For any dataset, we use 90\% of trials for model fitting (training set) and 10\% of trials for evaluating the quality of fit (held-out or test set).
The quality-of-fit statistic is negative log-likelihood (NLL) on held-out data, which is the sum of the negative log-likelihoods of all held-out trials.

Note that the model is fitted separately for each potential loss function exponent $p$ (giving a different heldout NLL for each $p$), whereas the other parameters (prior, encoding, noise and guessing parameters) are jointly fitted.


By Bayes' rule, the posterior over $\theta$ given an encoding $m$ in the $l$-th noise level ($l = 1, \dots, L$) is
\begin{equation}
P(\theta|m; \sigma_l) \propto P_{sensory}(m|\theta; \sigma_l^2) \Prior(\theta)
\end{equation}
where $P_{sensory}(m|\theta; \sigma_l^2)$ is the sensory noise distribution over $\mathcal{Y}$ (Gaussian with variance $\sigma_l^2$ or von Mises with concentration $1/\sigma_l^2$, for the corresponding noise level $l$).
The $L_p$-estimator $\hat{\theta}$ is obtained from the posterior $P(\theta|m; \sigma_l^2)$ by minimizing the $L_p$ loss function.
Then, the data log-likelihood for a single pair of true stimulus $\theta \in \mathcal{X}$ and observed response $x \in \mathcal{X}$ at the $l$-th noise level ($l=1, \dots, L$) is defined as follows:
\begin{equation}\label{eq:data-likelihood-raw}
\log P(x|\theta; \sigma_l^2) = \log \left\{ \frac{\xi}{\operatorname{Vol}(\mathcal{X})} + (1-\xi)  \int_{m \in \mathcal{Y}} P_{motor}(x|\hat{\theta}_m) P_{sensory}(m|\theta; \sigma_l^2) dm \right\}
\end{equation}
where $\hat{\theta}_m \in \mathcal{X}$ is the $L_p$-estimator given the posterior $P(\hat{\theta}|m; \sigma_l^2)$ for an encoding $m \in \mathcal{Y}$, and $P_{motor}(x|\hat{\theta}_m)$ is the motor noise distribution over $\mathcal{X}$ (Gaussian with variance $\rho^2$ or von Mises with concentration $1/\rho^2$).
This can equivalently be expressed, by transforming the integral over $m$ to an integral over $F^{-1}(m)$, as
\begin{equation}\label{eq:data-likelihood}
\log P(x|\theta; \sigma_l^2) = \log \left\{ \frac{\xi}{\operatorname{Vol}(\mathcal{X})} + (1-\xi)  \int_{m \in \mathcal{X}} P_{motor}(x|\hat{\theta}_m) P_{sensory-transformed}(m|\theta; \sigma_l^2) dm \right\}
\end{equation}
where, by the transformation law for probability densities,
\begin{equation}
P_{sensory-transformed}(m|\theta; \sigma_l^2) \propto P_{sensory}(F(m)|\theta; \sigma_l^2) \cdot F'(m)
\end{equation}
This is advantageous for implementation because it avoids the need to separately represent $\mathcal{X}$ and $\mathcal{Y}$.

\paragraph*{Parameterization of Prior and Encoding}
The model is implemented computationally on a regularly spaced discretized grid $\theta_1, \dots, \theta_N$.
We parameterize the encoding and prior by assigning one value of $F'$, $\Prior$ to each grid point $\theta_i$.
The encoding and prior are parameterized in terms of logarithms of these values, and obtained by normalizing the exponentiated parameters (softmax transform), as follows:
\begin{align}\label{eq:fprime-param}
	F'(\theta_i) =& \operatorname{Vol}(\mathcal{Y}) \frac{\exp(\alpha(\theta_i))}{\sum_{j=1}^N\exp(\alpha(\theta_j))} \\
\label{eq:prior-param}
		\Prior(\theta_i) =& \frac{\exp(\beta(\theta_i))}{\sum_{j=1}^N\exp(\beta(\theta_j))}
\end{align}
for $\theta_i$ on the grid; here, $\alpha(\theta_i)$ and $\beta(\theta_i)$ ($i=1,\dots, N$) are real-valued parameters of the model. 
We note that this is a discretized version of the theoretical construction of a measure over models (Section~\ref{sec:measure-mu}).

\paragraph*{Implementation of Model}
We map the  input stimulus $\theta$  to the closest grid point $\theta_i$ ($i = 1,\dots, N$) and then compute likelihoods and posteriors exactly 
on the discretized grid, and obtain the $L_p$ estimator via Newton's method.
First, the sensory likelihood is, for the $l$-th noise level ($k, i = 1, \dots, N$):
\begin{equation}
P_{sensory-transformed}(\theta_k|\theta_i; \sigma_l^2) \propto \exp(-\frac{(F(\theta_k) - F(\theta_i))^2}{2\sigma_l^2}) \cdot F'(\theta_k)
\end{equation}
(analogously in the case of von Mises noise).
Second, the posterior is ($j, k = 1, \dots, N$):
\begin{equation}
P(\theta_j|\theta_k; \sigma_l^2) \propto P_{sensory-transformed}(\theta_k|\theta_j; \sigma_l^2) \Prior(\theta_j)
\end{equation}
The estimator $\hat{\theta} \in \mathcal{X}$ is obtained from the posterior $P(\theta_k|\theta_i; \sigma_l^2)$ by minimizing the expected $L_p$ loss function:
\begin{equation}\label{eq:lp-estimator-discretized}
\hat{\theta} = \argmin_{\hat{\theta} \in \mathcal{X}} \sum_{k=1}^N P(\theta_k|\theta_i; \sigma_l^2) |\hat{\theta}-\theta_k|^p
\end{equation}
or, in the case of a circular space,
\begin{equation}\label{eq:lp-estimator-discretized-circular}
\hat{\theta} = \argmin_{\hat{\theta} \in \mathcal{X}} \sum_{k=1}^N P(\theta_k|\theta_i; \sigma_l^2) (1-\cos(\hat\theta-\theta))^{p/2}
\end{equation}
These formulas are used for $p \geq 2$; implementations for $p=0,1$ are described in Section~\ref{sec:implementation}. 
%
%On interval spaces, the $L_p$ loss function is defined as $|\hat\theta-\theta|^p$ (for $p>0$).
%On circular spaces, the $L_p$ loss function is defined as $(1-\cos(\hat\theta-\theta))^{p/2}$ (for $p>1$).
%The $L_0$ estimator is operationalized as the maximum of a kernel-smoothed version of the posterior (see \citet{hahn2024unifying} for details).

\paragraph*{Model Fitting}
The data log-likelihood for a single pair consisting of a true stimulus mapped to a grid point $\theta_i$ ($i=1, \dots, N$) and an observed response $x \in \mathcal{X}$, under the $l$-th noise level ($l=1,\dots,L$), is defined by discretizing the integral in Eq.~\eqref{eq:data-likelihood}:
\begin{equation}\label{eq:data-likelihood-discretized}
\log P(x|\theta_i; \sigma_l^2) = \log \left\{\frac{\xi}{\operatorname{Vol}(\mathcal{X})} + (1-\xi)  \sum_k P_{motor}(x|\hat{\theta}_{\theta_k}) P_{sensory-transformed}(\theta_k|\theta_i; \sigma_l^2) \right\}
\end{equation}
where $k$ runs over the points on the discretized grid, and $\hat{\theta}_{\theta_k}$ is the $L_p$ estimator derived from posterior $P(\theta|\theta_k; \sigma_l^2)$.
Summing this over all datapoints gives the log-likelihood of the behavioral dataset.

We fit all model parameters (prior, encoding, noise magnitudes, guessing rate) end-to-end via gradient descent to maximize the likelihood of the behavioral data.
We implement (\ref{eq:data-likelihood-discretized}) in PyTorch \citep{DBLP:conf/nips/PaszkeGMLBCKLGA19}, which provides differentiation out-of-the box.
Notably, we provide a custom implementation for the $L_p$ optimizers as standalone PyTorch modules, allowing differentiating through $\hat{\theta}$. In the forward pass, this module obtains $\hat{\theta}$ from the posterior $P(\theta|\theta_k; \sigma_l^2)$ by minimizing (\ref{eq:lp-estimator-discretized}--\ref{eq:lp-estimator-discretized-circular}) via Newton's method. For the backward pass, we compute the derivatives of $\hat{\theta}$ w.r.t. the weights of the posterior $P(\theta_j|\theta_k; \sigma_l^2)$ ($j, k = 1, \dots, N$) via implicit differentiation.
We refer to \citet{hahn2024unifying} for technical details.

\citet{hahn2024unifying} mitigate overfitting in the shapes of prior and encoding by adding a regularization term penalizing\footnote{For circular spaces, $\left(\alpha(\theta_N), \alpha(\theta_{1})\right)^2$ is added additionally.} $\sum_{i=1}^{N-1} \left(\alpha(\theta_i), \alpha(\theta_{i+1})\right)^2$ and similarly for $\beta$, modulated with a regularization strength $\lambda$ that is set by the user.
In principle, one can determine the regularization strength for a given behavioral dataset via cross-validation, though this comes with computational cost. 
As described in Materials and Methods, we used the same regularization strength as was identified for the data from \citet{deGardelle2010AnOI} by \citet{hahn2024unifying}, and used this across all datasets.
A control study (Section~\ref{sec:robustness-reg}) shows that the regularization strength has little impact on the identifiability.





\end{document}
